% Chapter 4: heat-generation extension of the graphite-anode ICA tail theory.
\documentclass[11pt,a4paper]{article}
\usepackage[margin=22mm]{geometry}
\usepackage{kotex}
\usepackage{amsmath,amssymb,mathtools,bm}
\usepackage{booktabs,longtable,array}
\usepackage{enumitem}
\usepackage{amsthm}
\usepackage{hyperref}
\usepackage{fancyhdr}
\usepackage{lastpage}
\usepackage{setspace}
\setstretch{1.13}
\setlength{\parskip}{0.45em}\setlength{\parindent}{0pt}\setlength{\headheight}{14pt}
\hypersetup{colorlinks=true, linkcolor=blue, urlcolor=blue, citecolor=blue,
  pdftitle={흑연 음극 ICA tail — Chapter 4}}
\pagestyle{fancy}\fancyhf{}
\lhead{흑연 음극 ICA tail — Ch.4 (entropy production 기반 발열 일반화)}\rhead{\thepage/\pageref{LastPage}}
\renewcommand{\headrulewidth}{0.4pt}

\newtheorem*{groundingbox}{Grounding}
\newtheorem*{boundbox}{유효범위(Validity bound)}
\newtheorem*{flagbox}{FLAGGED (미확정)}
\newtheorem*{keybox}{Keystone}
\newtheorem*{linkbox}{Ch1--3 전달식}
\newtheorem*{verifybox}{정합 검증 (미시=거시, P3-6)}
\newtheorem*{practicebox}{실무 팁 (본문 물리식 아님)}

\newcommand{\dd}{\mathrm{d}}
\newcommand{\eff}{\mathrm{eff}}
\newcommand{\eq}{\mathrm{eq}}
\newcommand{\app}{\mathrm{app}}
\newcommand{\drive}{\mathrm{drive}}
\newcommand{\bg}{\mathrm{bg}}
\newcommand{\cell}{\mathrm{cell}}
\newcommand{\ext}{\mathrm{ext}}
\newcommand{\OCV}{\mathrm{OCV}}
\newcommand{\tot}{\mathrm{tot}}
\newcommand{\irr}{\mathrm{irr}}
\newcommand{\rev}{\mathrm{rev}}
\newcommand{\rxn}{\mathrm{rxn}}
\newcommand{\ct}{\mathrm{ct}}
\newcommand{\ohm}{\mathrm{ohm}}
\newcommand{\conc}{\mathrm{conc}}
\newcommand{\src}{\mathrm{src}}
\newcommand{\loss}{\mathrm{loss}}
\newcommand{\rlx}{\mathrm{relax}}
\newcommand{\cand}{\mathrm{cand}}
\newcommand{\conv}{\mathrm{conv}}
\newcommand{\prog}{\mathrm{prog}}
\newcommand{\coord}{\mathrm{coord}}
\newcommand{\transp}{\mathrm{transport}}
\newcommand{\tr}{\mathrm{tr}}
\newcommand{\chem}{\mathrm{chem}}
\newcommand{\amb}{\mathrm{amb}}
\newcommand{\sstat}{\mathrm{ss}}
\newcommand{\AL}[1]{\textsf{[AL-#1]}}
\newcommand{\brk}{\discretionary{}{}{}}% 표/참고문헌 긴 DOI·ISBN 줄바꿈 허용(zero-width)

\begin{document}
\sloppy
\begin{center}
{\LARGE\bfseries 리튬이온전지 흑연 음극 ICA($dQ/dV$) tail 이론 — Chapter 4\par}
\vspace{0.4em}
{\large entropy-production 기반 발열 일반화: transition-level entropy production $\to$
비가역열 일반형 $\sum_j n_j^{\eff}I_j\eta_j$\par}
\vspace{0.2em}
{\normalsize Ch3 $r_j^\pm$ $\to$ $\dot{\mathcal Q}_{\irr,j}^{\tr}{=}\tfrac{Q_{j,\tot}}{F}RT(\mathcal J_j^+{-}\mathcal J_j^-)\ln\tfrac{\mathcal J_j^+}{\mathcal J_j^-}\ge0$ $\to$ 미시$=$거시 $n_j^{\eff}I_j\eta_j$\par}
\end{center}

% ====================================================================
\section*{서 — 인과 사슬에서 ``발열 정량화''의 자리}
\addcontentsline{toc}{section}{서 — 발열 정량화의 자리}
% ====================================================================
Chapter 1 은 ICA 꼬리의 주역을 \emph{동역학}(진행률 $\xi_j$ 의 유한속도 완화)으로 지목했고, Chapter 2 는 그
위에서 발열을 \emph{가역}(평형 entropy, 전위 온도계수)과 \emph{비가역}(과전압 소산, flux$\times$force) 기원으로
\emph{연결}했다. Chapter 3 은 그 비가역 소산의 동역학 원천인 mobility $k_j$ 를 \emph{미시} forward/backward
transition rate $r_j^\pm$(Level B)로 \emph{확대}하고, detailed balance 정합으로 effective electron number
$n_j^{\eff}=RT/(Fw_j)$ 를 고정했다. Chapter 4 는 이 셋을 합쳐 \textbf{하나의 정량 명제}를 닫는다.
\begin{quote}\itshape
Chapter 3 의 forward/backward rate $r_j^\pm$ 로부터 각 transition 의 \textbf{entropy production} 을 비평형
열역학의 표준형(network/master-equation)으로 세우면, 그것은 거시 비가역 발열 \emph{일반형}
$\dot{\mathcal Q}_\irr=\sum_j n_j^{\eff}I_j\eta_j\ge0$ 으로 \textbf{정확히 환원}된다. Chapter 2 가 쓴
$\sum_j I_j\eta_j$ 는 이 일반형의 ideal-width($n_j^{\eff}=1$) 특수해다.
\end{quote}
이 명제를 \textbf{다섯 단계}로 푼다. 각 단계는 학부 수준(비평형 열역학 flux$\times$force, Faraday 법칙,
지수$\cdot$로그 대수, 2법칙 $\dot S_\irr\ge0$)으로 따라갈 수 있게 중간을 생략하지 않는다.
\begin{enumerate}[label=\textbf{(\arabic*)},leftmargin=2.2em]
\item \textbf{비가역 열의 출발점은 flux$\times$force(entropy production)다.} 임의 heat correction 이 아니라
  conjugate flux $J_\alpha$ 와 thermodynamic force $X_\alpha$ 의 곱 $T\dot S_\irr=\sum_\alpha J_\alpha X_\alpha\ge0$
  이 2법칙이 보장하는 기본형이다(\S\ref{sec:ch4_irr}, \AL{40}).
\item \textbf{각 transition 의 entropy production 은 network 형으로 쓴다.} Chapter 3 의 forward/backward
  flux $\mathcal J_j^\pm$ 로 $\dot{\mathcal Q}_{\irr,j}^{\tr}=\tfrac{Q_{j,\tot}}{F}RT(\mathcal J_j^+-\mathcal J_j^-)
  \ln(\mathcal J_j^+/\mathcal J_j^-)\ge0$ 이 따라온다 — ``두 양수의 (차)$\times$(로그비)는 항상 음이 아니다''라는
  부등식이 2법칙을 직접 준다(\S\ref{sec:ch4_tr}, \AL{41},\AL{43}).
\item \textbf{미시 entropy production 이 거시 $n_j^{\eff}I_j\eta_j$ 로 환원된다.} detailed balance(Ch3)와
  $n_j^{\eff}=RT/(Fw_j)$(Ch3)를 대입하면 transition chemical affinity 가 $\mathcal A_j^{\chem}=n_j^{\eff}F\eta_j$
  로 정리되어 $\dot{\mathcal Q}_{\irr,j}^{\tr}=n_j^{\eff}I_j\eta_j$ 가 된다(\S\ref{sec:ch4_tr} verifybox,
  \AL{44}). 여기서 $\eta_j=V_{n,\drive}-V_{\eq,j}(\xi_j)$ 의 기준 전위를 \emph{단일}로 못박는다.
\item \textbf{일반형의 부호 양정치를 닫는다.} $\dot{\mathcal Q}_\irr=\sum_j n_j^{\eff}I_j\eta_j\ge0$ 이며
  $\sum_j I_j\eta_j$(Ch2)는 그 $n^{\eff}=1$ 특수해다(\S\ref{sec:ch4_general}, \AL{45}).
\item \textbf{가역 열$\cdot$transport$\cdot$휴지 relaxation 을 같은 원칙으로 정렬한다.} 가역 열은 평형 전위
  온도계수 또는 transition entropy basis(택일, 정합 제약), transport 는 flux$\times$force, 휴지 relaxation 은
  외부 $I=0$ 라도 내부 transition flux 의 dissipation 으로 남는다(\S\ref{sec:ch4_revOCV}--\ref{sec:ch4_rest}).
  Bernardi 가역열($q_\rev=s^\conv|I|T\,\partial U/\partial T$)이 미시 reaction entropy 의 합과 정합함을
  재확인한다.
\end{enumerate}

\noindent\textbf{한 문장 요지.} \emph{발열은 새 피팅항이 아니라, 이미 깔린 인과 사슬
($V_n\to\xi_{\eq,j}\to k_j/r_j^\pm\to\dd\xi_j/\dd t$) 위에서 계산되는 에너지 계층이며, 비가역 열은
transition-level entropy production 의 거시 표현 $\sum_j n_j^{\eff}I_j\eta_j\ge0$ 으로 닫힌다.} 본 장은
Chapter 1--3 과 같이 \textbf{방전(탈리튬화) 방향}($s_{\phi,j}=+1$)을 기본으로 한다. 충전 branch/hysteresis
heat 의 branch 의존 부호와 branch 비대칭 affinity 일반화는 Chapter 5 로 보낸다.

\begin{groundingbox}
\textbf{지배 표준(GS, Ch1--3 계승) 및 발열 근거 원칙.} \textbf{(GS-1)} 모든 식은 물리적 의미를 prose 로 먼저
말한 뒤 수식화하고 중간 단계를 생략하지 않는다(학부 무비약). \textbf{(GS-2)} 결론은 출판 수준 엄밀성.
\textbf{(GS-3)} 모든 가정은 통합 Assumption Ledger(\AL{\#}: 논문/교과서 근거 + 검증 DOI)를 인용하며, 근거가
없으면 FLAGGED 로 표기하고 established 로 쓰지 않는다. \textbf{(GS-4)} 임의 clip/cap/softplus/step-function
정의식과 근거 없는 수치 임계값을 물리 모델로 쓰지 않는다(속도 제한은 물리적 병목으로, log 발산은 domain guard
로). \textbf{발열 근거 원칙.} 모든 발열 항은 (i) entropy production, (ii) conjugate flux$\times$force, (iii)
평형 전위 온도계수, (iv) transition entropy$\times$진행률, (v) transport/finite-current dissipative loss, (vi)
자유에너지 감소를 동반하는 internal relaxation 중 하나의 \emph{물리} 근거를 가져야 한다. 임의 heat correction /
capped heat / empirical gain / residual heat 는 기본식에 넣지 않는다(실무 팁 분리). 방전 기준 $s_{\phi,j}=+1$,
$|I|>0$, heat-generation-positive convention.
\end{groundingbox}

\tableofcontents
\newpage

% ====================================================================
\section{Chapter 1--3 에서 전달받는 고정 기준식}\label{sec:ch4_inherit}
% ====================================================================
본 장은 Chapter 1$\cdot$2$\cdot$3 의 다음 기준식을 입력으로 받는다.
각 식 번호는 앞 장의 boxed 결과를 가리키며, 본 장은 그 위에 entropy-production 계층을 \emph{적층}한다(전하
보존식을 발열로 대체하지 않는다).
\begin{linkbox}
\textbf{(L1) 전하 보존식(중심 상태식, 발열로 대체 안 됨).} $Q_\cell\,q=Q_\bg^{\eq}(V_n,T)+\sum_{j=1}^{N_p}Q_{j,\tot}\xi_j$.
$V_n$ 은 그 해(외부 입력 아님); 평형 특수해가 $V_{n,\OCV}(q,T)$. 단조성 floor $\partial Q_\bg^{\eq}/\partial V_n\ge\epsilon_Q>0$
로 해 유일.\\
\textbf{(L2) kinetics(Ch3 미시 + Ch1 축약).} $\dfrac{\dd\xi_j}{\dd t}=r_j^+(1-\xi_j)-r_j^-\xi_j$, 축약형
$=k_j^{\mathrm{phys}}(\xi_{\eq,j}-\xi_j)$ ($k_j^{\mathrm{phys}}=r_j^++r_j^-$, mobility); 정전류 구간서만
$\dfrac{\dd\xi_j}{\dd t}=\dfrac{|I|}{Q_\cell}\dfrac{\dd\xi_j}{\dd q}$.\\
\textbf{(L3) Ch2 dissipation force(비가역 열 공액 force).} $\eta_j=V_{n,\drive}-V_{\eq,j}(\xi_j)$,
$V_{\eq,j}(\xi_j)=U_j(T)+w_j(T)\ln[\xi_j/(1-\xi_j)]$ — 현재 $\xi_j$ 의 \emph{국소 평형} 기준(중심 affinity
$\mathcal A_j$ 와 구분). \textbf{기준 전위는 $V_{n,\drive}$ 단일}(\S\ref{sec:ch4_eta_single}).\\
\textbf{(L4) Ch3 미시 정합.} detailed balance $r_j^+/r_j^-=\xi_{\eq,j}/(1-\xi_{\eq,j})$, 그 ln 형
$\ln[\xi_{\eq,j}/(1-\xi_{\eq,j})]=(V_n-U_j)/w_j$; effective electron number $n_j^{\eff}=RT/(Fw_j)$ ($\mathcal A_j
=s_{\phi,j}n_j^{\eff}F(V_{n,\drive}-U_j)$, $s_{\phi,j}=+1$ 방전 기본). Chapter 1 의 $\chi_j$ 는 Level-A scalar
mobility sensitivity 이고, Chapter 3 의 $\beta_j$ 는 Level-B directional transfer coefficient 이므로 동일물로
두지 않는다. 둘의 비교는 Chapter 3 의 projection $\chi_{j,\mathrm{proj}}$ 를 통해서만 한다.\\
\textbf{(L5) 네 전위 구분.} 내부 $V_n$(보존식 해), apparent $V_{n,\app}=V_n+s_I|I|R_n$(분극 포함, 평형 전위
아님), 구동 $V_{n,\drive}=V_n+\eta_{n,\drive}$(기본 $\eta_{n,\drive}=0\Rightarrow V_{n,\drive}=V_n$), 평형
$V_{n,\OCV}$(평형 보존식 해).\\
\textbf{(L6) Ch2 가역열 경계.} $q_{\rev,j}=I_j\,T\,\partial U_j/\partial T=I_j\Pi_j$ ($\Pi_j=T\partial U_j/\partial T$,
Peltier-type); reduced $\dot{\mathcal Q}_{\rev,n}^{\OCV}=s_{\rev,n}^{\conv}|I|T(\partial V_{n,\OCV}/\partial T)_q$;
$\partial U_j/\partial T=\Delta S_{\rxn,j}/(s_{\phi,j}F)$ (reaction entropy, activation entropy $\Delta S_{a,j}$
와 별개).\\
\textbf{(L7) Ch2 비가역열 특수형($n^{\eff}=1$).} flux$\times$force $T\dot S_{\irr,n}=\sum_j J_{n,j}F\eta_j\ge0$,
축약 $q_\irr=|I|\eta\ge0$. \textbf{본 장이 일반화하는 대상} — Ch2 가 forward-ref 한 일반형 $\sum_j n_j^{\eff}I_j\eta_j$
를 본 장에서 닫는다(\S\ref{sec:ch4_general}).
\end{linkbox}
임의 heat correction / capped heat / empirical residual 을 기본 발열식으로 쓰지 않는다(GS-4); 강구동 제한과
log 발산은 물리적 병목과 numerical domain guard 로 둔다.

\textbf{신규 기호(Ch1--3 위에 추가).} 본 장에서만 새로 도입하는 기호를 모은다(앞 장 기호는 (L1)--(L7) 과 동일물).
\begin{longtable}{p{0.16\textwidth} p{0.13\textwidth} >{\raggedright\arraybackslash}p{0.63\textwidth}}
\toprule 기호 & 단위 & 의미 \\ \midrule \endhead
$\dot S_\irr$ & W/K & entropy production rate (비가역 entropy 생성률; 2법칙 $\ge0$) \\
$J_\alpha,\ X_\alpha$ & (쌍) & generalized conjugate flux $\cdot$ thermodynamic force (곱 $J_\alpha X_\alpha$ 가 W) \\
$J_n,\ \mathcal A_n$ & mol/s, J/mol & 음극 molar reaction flux $\cdot$ generalized reaction affinity \\
$\mathcal J_j^+,\ \mathcal J_j^-$ & 1/s & transition $j$ 의 forward $=r_j^+(1-\xi_j)$ $\cdot$ backward $=r_j^-\xi_j$ unidirectional flux \\
$\mathcal A_j^{\chem}$ & J/mol & transition $j$ 의 chemical affinity $=RT\ln(\mathcal J_j^+/\mathcal J_j^-)$ (network thermo) \\
$I_j$ & A & transition $j$ lumped current $=Q_{j,\tot}\,\dd\xi_j/\dd t$ (Ch3 동일물) \\
$\dot n_j$ & mol/s & transition $j$ molar rate $=(Q_{j,\tot}/F)\,\dd\xi_j/\dd t$ \\
$\dot{\mathcal Q}_{\irr,j}^{\tr}$ & W & transition $j$ 의 transition-level 비가역 발열률 (entropy production $\times T$) \\
$\dot{\mathcal Q}_{\src,n}$ & W & 음극 control volume 내부 총 열원 (가역 흡열 가능) \\
$\dot{\mathcal Q}_{\irr,n},\dot{\mathcal Q}_{\rev,n}$ & W & 계면 반응 비가역 발열 $\cdot$ 가역 발열 (음극 합) \\
$\dot{\mathcal Q}_{\rlx,n},\dot{\mathcal Q}_{\transp,n}$ & W & 내부 relaxation 발열 $\cdot$ transport/병목 dissipative loss \\
$\Delta S_j(\equiv\Delta S_{\rxn,j})$ & J/(mol K) & transition $j$ 의 reaction entropy (1 mol electron/Li-equiv 당) \\
$h_\bg$ & V/K & background entropy coefficient $=(\partial V_{\bg,\eq}/\partial T)$ (Ch2 동일물) \\
$R_{\transp,n},R_{n,\eff}$ & $\Omega$ & transport effective resistance $\cdot$ heat-generation effective resistance \\
$J_\ell,\ X_\ell$ & (쌍) & transport flux $\cdot$ 그 공액 force (chemical/electrolyte-potential/concentration gradient) \\
$C_{\mathrm{th}}$ & J/K & 음극 control volume 유효 열용량 (Ch2 동일물) \\
$s^{\conv},\ s_\bg^\conv$ & --- & convention bookkeeping factor $\in\{+1,-1\}$ (피팅 X; Ch2 동일물) \\
\bottomrule
\end{longtable}
$F=96485\ \mathrm{C/mol}$(Faraday), $R=8.314\ \mathrm{J/(mol\,K)}$(기체상수), $\mathrm{J/C}=\mathrm V$. 열률 dot
표기는 시간율 [W]. unidirectional flux $\mathcal J_j^\pm$ 와 net flux $\dd\xi_j/\dd t=\mathcal J_j^+-\mathcal J_j^-$
의 단위는 $\mathrm{s^{-1}}$($\xi_j$ 무차원).

\textbf{직접 입력은 시간 미분량.} 모든 발열률$\cdot$entropy production 의 직접 입력은 시간 미분량 $\dd\xi_j/\dd t$
또는 시간당 flux 다. transition 별 전하/molar 진행률은(Ch3 \eqref{eq:ch4_Ij} 와 동일 정의)
\begin{equation}
I_j=Q_{j,\tot}\,\frac{\dd\xi_j}{\dd t},\qquad
\dot n_j=\frac{Q_{j,\tot}}{F}\,\frac{\dd\xi_j}{\dd t}\ [\mathrm{mol\,s^{-1}}].
\label{eq:ch4_Ij}
\end{equation}
\textbf{차원 검산.} $I_j=Q_{j,\tot}[\mathrm C]\cdot\dd\xi_j/\dd t[\mathrm{s^{-1}}]=\mathrm{C/s}=\mathrm A$;
$\dot n_j=(Q_{j,\tot}/F)[\mathrm{C}/(\mathrm{C\,mol^{-1}})]\cdot[\mathrm{s^{-1}}]=\mathrm{mol/s}$ — 정합. 면적
정규화가 필요하면 유효 반응면적 $A_{\eff}$ 로 area-flux $J_j=\dot n_j/A_{\eff}$ 를 \emph{별도} 정의한다(lumped
molar rate $\dot n_j$ 와 area-flux $J_j$ 혼동 금지).

% ====================================================================
\section{$\eta_j$ 의 기준 전위 단일화}\label{sec:ch4_eta_single}
% ====================================================================
본 절은 \textbf{본 장의 가장 먼저 못박아야 할 규약}이다. Chapter 2 는 dissipation force(과전압)를
$\eta_j=V_{n,\drive}-V_{\eq,j}(\xi_j)$ 로(구동 전위 기준) 정의했고, Chapter 3 은 중심 기준 affinity 의
전위 단위 형 $\eta_j=s_{\phi,j}(V_{n,\drive}-U_j)$ 를 별도로 썼다. 본 장의 \emph{발열} 유도(특히 미시=거시
정합, \S\ref{sec:ch4_tr})에서는 ``비가역 열 공액 force''를 쓰므로, 그것을 \textbf{Chapter 2 의 dissipation
force 하나로 단일화}한다(\AL{42}).
\begin{keybox}
\textbf{본 장 비가역 열 공액 force 의 단일 정의.} 비가역 발열의 공액 force 는 \emph{현재}
$\xi_j$ 의 국소 평형으로부터의 이탈, 즉
\begin{equation}
\boxed{\;\eta_j\;\equiv\;V_{n,\drive}-V_{\eq,j}(\xi_j)\;,\qquad V_{\eq,j}(\xi_j)=U_j(T)+w_j(T)\ln\frac{\xi_j}{1-\xi_j}\;,}
\label{eq:ch4_eta_def}
\end{equation}
이며 \textbf{기준 전위는 $V_{n,\drive}$ 단일}이다($V_n$ 과 $V_{n,\drive}$ 혼용 0). 기본 전개에서는
\textbf{기본형 $V_{n,\drive}=V_n$}($\eta_{n,\drive}=0$, 식 (L5))을 쓰며, 이 전제는 미시=거시 대입
(\S\ref{sec:ch4_tr})에서 \emph{명시적으로} 사용된다. 강구동($\eta_{n,\drive}\ne0$, 따라서 $V_{n,\drive}\ne V_n$)
일반화는 BOUNDED 로 표기하고 유효범위를 병기한다.
\end{keybox}
\textbf{왜 단일화가 필요한가(이중정의의 위험).} 만약 $\eta_j$ 를 어떤 식에서는 $V_n$ 기준
($V_n-V_{\eq,j}(\xi_j)$), 다른 식에서는 $V_{n,\drive}$ 기준($V_{n,\drive}-V_{\eq,j}(\xi_j)$)으로 섞어 쓰면,
$V_{n,\drive}\ne V_n$ 인 강구동에서 \emph{같은 기호 $\eta_j$ 가 서로 다른 값}을 가리켜 비가역 발열의 부호와
크기가 식마다 달라진다 — 이것이 Phase A 가 CRITICAL 로 적발한 ``$\eta$ 이중정의''다. 본 장은 \emph{정의}
\eqref{eq:ch4_eta_def} 를 $V_{n,\drive}$ 기준으로 동결하고, 기본형에서 $V_{n,\drive}=V_n$ 임을 사용 시점마다
명시해 모순을 차단한다(Chapter 2 \S 비가역 열 부호 정합 + Chapter 3 \S$V_{n,\drive}$ 정의장과 동일 입도).

\textbf{$V_{\eq,j}(\xi_j)$ 의 유도(logistic 역함수; Ch2 재확인, 무비약).} Chapter 1 logistic(식 (L4),
$s_{\phi,j}=+1$) $\xi_{\eq,j}=[1+e^{-(V_n-U_j)/w_j}]^{-1}$ 을 ``$\xi_j$ 가 현재값일 때 그것을 평형으로 만드는
전위''에 대해 역으로 푼다. $\xi_j=[1+e^{-(V_{\eq,j}-U_j)/w_j}]^{-1}$ 로 두고 정리하면 $1/\xi_j-1=
e^{-(V_{\eq,j}-U_j)/w_j}$, 즉 $(1-\xi_j)/\xi_j=e^{-(V_{\eq,j}-U_j)/w_j}$. 양변에 자연로그를 취하면
$\ln[(1-\xi_j)/\xi_j]=-(V_{\eq,j}-U_j)/w_j$, 부호를 뒤집어 $\ln[\xi_j/(1-\xi_j)]=(V_{\eq,j}-U_j)/w_j$, 따라서
$V_{\eq,j}(\xi_j)=U_j+w_j\ln[\xi_j/(1-\xi_j)]$ — 식~\eqref{eq:ch4_eta_def} 의 둘째 식이다(중간 단계 전부 명시;
Chapter 2 의 동일 유도를 본 장 기준 전위 단일화 위에서 재확인).
\begin{boundbox}
\textbf{rate-affinity $\mathcal A_j$ 와 heat-force $\eta_j$ 의 구분(\AL{42}; Ch2$\cdot$Ch3 계승).} 식 (L4) 의
$\mathcal A_j=s_{\phi,j}n_j^{\eff}F(V_{n,\drive}-U_j)$ 는 \emph{중심} $U_j$ 기준이라 평형에서도 0 이 아니다
(rate 구동용). 본 절 $\eta_j$(식~\eqref{eq:ch4_eta_def})는 \emph{현재} $\xi_j$ 의 국소 평형 $V_{\eq,j}(\xi_j)$
기준이라 평형($\xi_j=\xi_{\eq,j}$)에서 0 이 된다(dissipation 용). 둘의 차이는 식~\eqref{eq:ch4_eta_def} 를
$\mathcal A_j$ 의 전위 단위 형 $\mathcal A_j/(n_j^{\eff}F)=s_{\phi,j}(V_{n,\drive}-U_j)$ 와 비교하면(방전
$s_{\phi,j}=+1$, 기본형 $V_{n,\drive}=V_n$)
\[
\frac{\mathcal A_j}{n_j^{\eff}F}=(V_{n,\drive}-U_j)
=\underbrace{(V_{n,\drive}-V_{\eq,j}(\xi_j))}_{=\eta_j}
+\underbrace{(V_{\eq,j}(\xi_j)-U_j)}_{=w_j\ln[\xi_j/(1-\xi_j)]},
\]
즉 $\mathcal A_j/(n_j^{\eff}F)=\eta_j+w_j\ln[\xi_j/(1-\xi_j)]$ 이며, 둘째 항(configurational 가역 혼합 entropy)은
\emph{가역} entropy basis(\S\ref{sec:ch4_revxi})로 가고 \emph{비가역} dissipation 에는 $\eta_j$ 만 쓴다 — 이
구분을 혼동하면 가역 혼합 entropy 가 비가역 열로 오계상된다(BOUNDED — logistic 평형 모델 내에서 정확).
\end{boundbox}

% ====================================================================
\section{내부 열원 항의 물리적 분해 (Ch2 3-분해의 정련)}\label{sec:ch4_decomp}
% ====================================================================
\textbf{물리.} 음극 control volume 내부 열원을 \emph{물리 기원별}로 분해한다(피팅 편의 분해가 아니라
entropy-production / 평형 entropy / dissipative-loss 기원 분해; \cite{bernardi1985,rao1997}, \AL{46}):
\begin{equation}
\dot{\mathcal Q}_{\src,n}=\dot{\mathcal Q}_{\irr,n}+\dot{\mathcal Q}_{\rev,n}+\dot{\mathcal Q}_{\rlx,n}+\dot{\mathcal Q}_{\transp,n}.
\label{eq:ch4_decomp}
\end{equation}
\begin{longtable}{p{0.28\textwidth} >{\raggedright\arraybackslash}p{0.60\textwidth}}
\toprule 항 & 물리적 의미 \\ \midrule \endhead
$\dot{\mathcal Q}_{\irr,n}$ & 계면 반응/transition 진행의 비가역 entropy production ($\sum_j n_j^{\eff}I_j\eta_j$, $\ge0$) \\
$\dot{\mathcal Q}_{\rev,n}$ & 평형 entropy coefficient 또는 transition entropy 가역 열 (열역학 기원, 부호가변) \\
$\dot{\mathcal Q}_{\rlx,n}$ & 외부 전류와 무관하게 내부 상태가 자유에너지를 낮추며 relax 할 때의 열 (확정엔 자유에너지 함수 필요) \\
$\dot{\mathcal Q}_{\transp,n}$ & solid diffusion/electrolyte/film/finite-current limitation 의 dissipative loss ($\ge0$) \\
\bottomrule
\end{longtable}
\textbf{Chapter 2 3-분해와의 관계(정련, 이중계상 금지; \AL{46}).} 본 4-분해는 Chapter 2 의 3-분해
$\dot{\mathcal Q}_{\src,n}=\dot{\mathcal Q}_{\irr,n}+\dot{\mathcal Q}_{\rev,n}+\dot{\mathcal Q}_{\rlx,n}^{\cand}$
를 \emph{정련}한 것이다 — Chapter 2 에서 $\dot{\mathcal Q}_{\irr,n}$ 에 묶여 있던 transport/diffusion
dissipation 을 본 장에서 $\dot{\mathcal Q}_{\transp,n}$ 으로 \emph{분리}했고, 후보로만 남았던 relaxation heat
$\dot{\mathcal Q}_{\rlx,n}^{\cand}$ 을 본 장에서 entropy production 으로 \emph{확정 분해}한다(\S\ref{sec:ch4_rest}).
즉 기호 관계는
\begin{equation}
\dot{\mathcal Q}_{\irr,n}^{\mathrm{Ch2}}\;\supset\;\dot{\mathcal Q}_{\irr,n}^{\mathrm{Ch4}}+\dot{\mathcal Q}_{\transp,n},
\label{eq:ch4_ch2refine}
\end{equation}
이며, 두 장의 항을 함께 쓸 때 transport dissipation 을 \emph{이중 계상하지 않는다}(본 장 $\dot{\mathcal Q}_{\irr,n}$
은 계면 반응/transition entropy production 만 포함). 가역 열이 흡열일 수 있으므로 $\dot{\mathcal Q}_{\src,n}$
은 항상 양수 ``생성량''이 아니라 internal heat source term 이다. 외부 손실 포함 control-volume 1법칙 열수지는
(Chapter 2 식과 동형)
\begin{equation}
C_{\mathrm{th}}\frac{\dd T}{\dd t}=\dot{\mathcal Q}_{\src,n}-\dot{\mathcal Q}_{\loss},\qquad
\dot{\mathcal Q}_{\loss}=hA(T-T_\amb),
\label{eq:ch4_balance}
\end{equation}
(좌변 $\mathrm{(J/K)(K/s)=W}$ — 차원 정합). $\dot{\mathcal Q}_{\loss}$ 의 상세 열전달 모델(대류계수 $h$, 표면적
$A$, 주위 온도 $T_\amb$)은 본 장에서 닫지 않는다(Chapter 2 의 control-volume 형 계승).

% ====================================================================
\section{비가역 열의 출발점: flux--force 와 entropy production}\label{sec:ch4_irr}
% ====================================================================
\textbf{물리.} 비가역 열은 임의 heat correction 이 아니라 \emph{비평형 열역학의 entropy production} 에서 나온다.
가장 기본적 출발점은 conjugate flux $J_\alpha$ 와 그에 공액인 thermodynamic force $X_\alpha$ 의 곱이며, 2법칙이
그 합의 비음성을 보장한다(\cite{degrootmazur1962,onsager1931,prigogine1967}, \AL{40}):
\begin{equation}
\boxed{\;T\dot S_{\irr}=\sum_\alpha J_\alpha X_\alpha\ \ge\ 0\;.}
\label{eq:ch4_fluxforce}
\end{equation}
\textbf{2법칙 의미(무비약).} $\dot S_\irr$ 는 비가역 과정이 생성하는 entropy 율이며, 고립계에서 entropy 가
감소할 수 없다는 2법칙이 곧 $\dot S_\irr\ge0$ 이다. 각 $J_\alpha X_\alpha$ 가 \emph{같은 방향}으로 정렬된
공액 쌍(force 가 flux 를 그 방향으로 밀어냄)이면 곱이 비음이고, 그 합이 $T\dot S_\irr$ 다. 좌변에 $T$ 가 곱해진
것은 entropy production $\dot S_\irr$[W/K]을 \emph{열률}[W]로 환산했기 때문이다($T\dot S_\irr$ 가 dissipated
power). \textbf{차원 검산.} $J_\alpha X_\alpha$ 의 곱이 [W]이려면 force 와 flux 가 공액(예: molar flux
$[\mathrm{mol/s}]\times$ molar affinity $[\mathrm{J/mol}]=[\mathrm{J/s}]=\mathrm W$)이어야 한다 — 이 공액성이
flux$\times$force 의 정의 자체에 들어 있다.

\textbf{전기화학 반응의 specialization.} 음극 반응을 molar reaction flux $J_n$ [mol/s]와 generalized reaction
affinity $\mathcal A_n$ [J/mol]로 쓰면 식~\eqref{eq:ch4_fluxforce} 의 한 쌍이
\begin{equation}
\boxed{\;T\dot S_{\irr,n}=J_n\,\mathcal A_n\ \ge\ 0\;.}
\label{eq:ch4_JA}
\end{equation}
\textbf{$J_n\mathcal A_n$ 이 $I_n\eta_n$ 보다 기본형인 이유.} $I_n=FJ_n$, $\mathcal A_n=F\eta_n$ convention 이
일관되면 $J_n\mathcal A_n=I_n\eta_n$ 형과 연결되나, ``$I_n\eta_n$''의 부호는 전류 방향$\cdot$전극 전위
convention 에 종속된다(Chapter 2 가 $q_\irr=|I|\eta$ 로 부호를 양정치화한 이유와 같다). 따라서 본 장은
\emph{positive-dissipation} 형 식~\eqref{eq:ch4_JA} 를 더 기본으로 두고, magnitude 표현
$\dot{\mathcal Q}_{\irr,n}^{\mathrm{mag}}=|I_n||\eta_n|$ 은 \emph{부호가 정렬됐을 때의 크기} 표현이지 entropy
production 의 일반 정의가 아님을 명시한다(\AL{42}). 이로써 ``$I\eta$ 단독''(부호 미정렬)을 비가역 열로 쓰는
것을 차단한다($q_\irr=|I|\eta\ge0$ 형).

% ====================================================================
\section{Transition-level entropy production (미시) 과 거시식의 정합}\label{sec:ch4_tr}
% ====================================================================
이 절이 본 장의 핵심이다 — \emph{Chapter 3 의 forward/backward rate $r_j^\pm$ 가 어떻게 거시 비가역 발열
$n_j^{\eff}I_j\eta_j$ 로 환원되는가}.

\textbf{forward/backward unidirectional flux.} Chapter 3 transition kinetics(식 (L2))의 net flux 를 forward 와
backward 의 \emph{unidirectional} 성분으로 분해한다(\AL{41}):
\begin{equation}
\mathcal J_j^+=r_j^+\,(1-\xi_j),\qquad \mathcal J_j^-=r_j^-\,\xi_j,\qquad
\frac{\dd\xi_j}{\dd t}=\mathcal J_j^+-\mathcal J_j^-.
\label{eq:ch4_Jpm}
\end{equation}
$\mathcal J_j^+$ 는 미점유$\to$점유(방전 방향) 단위시간 진행, $\mathcal J_j^-$ 는 그 역방향이다(Chapter 3
식 $\dd\xi_j/\dd t=r_j^+(1-\xi_j)-r_j^-\xi_j$ 와 동일물; 본 장은 두 성분을 \emph{따로} 보존한다). \textbf{차원
검산.} $\mathcal J_j^\pm=r_j^\pm[\mathrm{s^{-1}}]\times$무차원 fraction $=\mathrm{s^{-1}}$, net 도 $\mathrm{s^{-1}}$
— 정합. 두 성분 모두 비음($r_j^\pm\ge0$, $0\le\xi_j\le1$).

\subsection{transition entropy production 의 network 형 (무비약)}\label{sec:ch4_trep}
\textbf{물리.} 비평형 열역학의 network/master-equation 이론에서, 두 상태 사이를 오가는 한 transition 의 entropy
production 은 \emph{net flux}와 \emph{forward/backward flux 비의 로그}의 곱이다(\cite{schnakenberg1976,prigogine1967},
\AL{41},\AL{43}). 이 ``$(\text{net flux})\times\ln(\text{flux 비})$'' 구조가 master equation 의 국소 detailed
balance 에서 따라오며, transition $j$ 의 진행이 동반하는 비가역 발열률은(몰 함량 $Q_{j,\tot}/F$ 와 $RT$ 로
extensive 화)
\begin{equation}
\boxed{\;\dot{\mathcal Q}_{\irr,j}^{\tr}=\frac{Q_{j,\tot}}{F}\,RT\,\big(\mathcal J_j^+-\mathcal J_j^-\big)\,
\ln\!\frac{\mathcal J_j^+}{\mathcal J_j^-}\ \ge\ 0\;.}
\label{eq:ch4_tr_entropy}
\end{equation}
\textbf{$\ge0$ 의 무비약 증명(부호 보조정리; 2법칙).} $\mathcal J_j^+,\mathcal J_j^->0$ 인 두 양수에 대해
$(\mathcal J_j^+-\mathcal J_j^-)\ln(\mathcal J_j^+/\mathcal J_j^-)\ge0$ 임을 직접 보인다. (i) $\mathcal J_j^+>\mathcal J_j^-$
이면 차 $>0$ 이고 비 $>1$ 이라 $\ln>0$ — 곱 $>0$. (ii) $\mathcal J_j^+<\mathcal J_j^-$ 이면 차 $<0$ 이고 비
$<1$ 이라 $\ln<0$ — 음$\times$음 $=$ 곱 $>0$. (iii) $\mathcal J_j^+=\mathcal J_j^-$(평형, net flux $0$)이면
차 $=0$ 이라 곱 $=0$. 세 경우 모두 $\ge0$ 이며, 이것이 식~\eqref{eq:ch4_tr_entropy} 의 2법칙 정합이다(``두
양수의 (차)$\times$(로그비)는 항상 음이 아니다''). \textbf{차원 검산.} $(Q_{j,\tot}/F)[\mathrm{mol}]\times
RT[\mathrm{J/mol}]\times(\mathcal J_j^+-\mathcal J_j^-)[\mathrm{s^{-1}}]\times\ln(\cdot)[\text{무차원}]
=\mathrm{mol\cdot(J/mol)\cdot s^{-1}}=\mathrm{J/s}=\mathrm W$ — 발열률의 올바른 단위. ($RT$ 가 $\mathrm{J/mol}$,
$Q_{j,\tot}/F$ 가 $\mathrm{mol}$ 이라 두 몰 인자가 곱해져 extensive [W]가 됨 — Chapter 1 $\rho_G$[mol/J]$\cdot$
Chapter 2 $\sigma_j$[W/mol] 의 ``intensive$\leftrightarrow$extensive 몰 인자 정합'' 교훈 계승.)

\textbf{extensive 화의 다리(무비약; $k_B\!\to\!R$, $\tfrac12$ 부재 사유).} Schnakenberg 의 단일 directed
transition(2-state site)에 대한 per-site entropy production 은 $\dot s_j^{\mathrm{site}}=k_B(\mathcal J_j^+-\mathcal J_j^-)
\ln(\mathcal J_j^+/\mathcal J_j^-)$ [W/K]다 — transition $j$ 를 \emph{단일 directed edge} 로 한 번만 세므로 양방향
edge 이중합의 $\tfrac12$ 인자는 \emph{없다}(이중합 $\tfrac12\sum_{\text{edges}}$ 와 단일-쌍 $\sum_j$ 는 동일물의
다른 셈법; 본 장은 후자). 여기에 (a) $T$ 를 곱해 per-site 발열률 [W], (b) site 수 $N_{\mathrm{site},j}=N_A(Q_{j,\tot}/F)$
(몰 함량 $\times$ Avogadro)를 곱해 독립 site 합, (c) $N_A k_B=R$ 로 묶으면 $N_A(Q_{j,\tot}/F)\,k_B T=(Q_{j,\tot}/F)RT$
가 되어 식~\eqref{eq:ch4_tr_entropy} 의 prefactor 가 무유도 점프 없이 따라온다(per-site $k_B$ $\to$ per-mole $R$
치환이 site 수 몰 인자와 정확히 짝).

\begin{verifybox}
\textbf{미시 entropy production $=$ 거시 $n_j^{\eff}I_j\eta_j$ (Ch2$\cdot$Ch3 정합, 무비약).}
transition chemical affinity 를 network thermo 정의 $\mathcal A_j^{\chem}=RT\ln(\mathcal J_j^+/\mathcal J_j^-)$
(\AL{43})로 두고, 그것을 Chapter 3 의 미시 정합으로 전개한다.

\emph{(1단계) flux 비를 rate 비$\times$점유 odds 로 분해.} 식~\eqref{eq:ch4_Jpm} 에서
$\dfrac{\mathcal J_j^+}{\mathcal J_j^-}=\dfrac{r_j^+(1-\xi_j)}{r_j^-\xi_j}=\dfrac{r_j^+}{r_j^-}\cdot\dfrac{1-\xi_j}{\xi_j}$.
따라서 $\ln(\mathcal J_j^+/\mathcal J_j^-)=\ln(r_j^+/r_j^-)+\ln[(1-\xi_j)/\xi_j]=\ln(r_j^+/r_j^-)-\ln[\xi_j/(1-\xi_j)]$.

\emph{(2단계) Chapter 3 detailed balance 대입(★ 두 전제 명시).} Chapter 3 의 detailed balance ln 형
$\ln(r_j^+/r_j^-)=(V_n-U_j)/w_j$ 를 넣는다. \textbf{이 ln 형은 두 전제에서만 닫힌형으로 성립한다}(Chapter 3
verifybox 한정과 \emph{동일 입도}): (i) ratio $r_j^+/r_j^-$ 가 $\xi_j$\emph{-무관}(즉 $C_j\equiv0$; Chapter 3 의
대칭 split, \AL{34})이라야 \emph{평형} odds 관계가 임의 \emph{비평형} $\xi_j$ 로 외삽되어 닫힌형이 되고, (ii)
\emph{내부 전위} $V_n$ 기준이다(detailed balance 는 평형 전하보존식의 해 $V_n$ 에서 성립). $C_j\not\equiv0$($r^+/r^-$
의 $\xi_j$-의존)이면 우변에 $RT\,C_j(\xi_j)$ 항이 추가되어 아래 닫힘이 깨지므로, 그 경우 미시=거시 환원은 BOUNDED 다
(\AL{44}; 본 장 기본 전개는 $C_j\equiv0$ 대칭 split). 전제 (i),(ii) 하에서
\begin{equation}
\mathcal A_j^{\chem}=RT\,\ln\!\frac{\mathcal J_j^+}{\mathcal J_j^-}
=RT\Big[\frac{V_n-U_j}{w_j}-\ln\frac{\xi_j}{1-\xi_j}\Big]
=\frac{RT}{w_j}\Big[(V_n-U_j)-w_j\ln\frac{\xi_j}{1-\xi_j}\Big].
\label{eq:ch4_affinity_raw}
\end{equation}

\emph{(3단계) 대괄호 $=$ overpotential $\eta_j$ (기본형 $V_{n,\drive}=V_n$ 전제 명시).}
대괄호 안은 $(V_n-U_j)-w_j\ln[\xi_j/(1-\xi_j)]$ 다. 여기서 \S\ref{sec:ch4_eta_single} 의 $V_{\eq,j}(\xi_j)=U_j
+w_j\ln[\xi_j/(1-\xi_j)]$ 를 빼는 형으로 묶으면 $(V_n-U_j)-w_j\ln[\xi_j/(1-\xi_j)]=V_n-V_{\eq,j}(\xi_j)$. 본 장
규약(식~\eqref{eq:ch4_eta_def})의 $\eta_j=V_{n,\drive}-V_{\eq,j}(\xi_j)$ 와 일치시키려면 \textbf{기본형
$V_{n,\drive}=V_n$}($\eta_{n,\drive}=0$)을 \emph{여기서 명시적으로 전제}한다 — 그래야
detailed balance 가 쓴 $V_n$ 과 dissipation force 가 쓴 $V_{n,\drive}$ 가 같은 무대에서 닫힌다. 이 전제 하에서
$(V_n-U_j)-w_j\ln[\xi_j/(1-\xi_j)]=V_{n,\drive}-V_{\eq,j}(\xi_j)=\eta_j$. (강구동 $V_{n,\drive}\ne V_n$ 에서는
detailed balance 의 $V_n$ 과 force 의 $V_{n,\drive}$ 가 갈리므로 이 닫힘이 일반적으로 성립하지 않는다 —
BOUNDED, \S\ref{sec:ch4_general}.)

\emph{(4단계) $n_j^{\eff}$ 흡수 + 부호 인자 $s_{\phi,j}$ 복원.} $RT/w_j$ 에 Chapter 3 의
$n_j^{\eff}=RT/(Fw_j)$(식 (L4))를 적용하면 $RT/w_j=n_j^{\eff}F$ 다. \emph{위 1$\sim$3단계 대수는 $s_{\phi,j}$ 를
포함하지 않으므로}, boxed 결과의 일반 $s_{\phi,j}$ 인자는 Chapter 3 명시형과의 \emph{표기 일관}을 위한 것이며,
충전 branch($s_{\phi,j}=-1$)에서 logistic 지수 부호로부터의 \emph{유도}는 Chapter 5 로 이월한다(본 장 유도는 방전
$s_{\phi,j}=+1$ 한정). 이 전제 하에서 식~\eqref{eq:ch4_affinity_raw} 는
\begin{equation}
\boxed{\;\mathcal A_j^{\chem}=\frac{RT}{w_j}\,\eta_j=n_j^{\eff}F\,\eta_j\;,\qquad
\eta_j=s_{\phi,j}\big[V_{n,\drive}-V_{\eq,j}(\xi_j)\big]\ \ (s_{\phi,j}=+1\ \text{방전}).}
\label{eq:ch4_affinity_eta}
\end{equation}
이는 Chapter 3 의 명시형 $\mathcal A_j=s_{\phi,j}n_j^{\eff}F(V_{n,\drive}-U_j)$ 와 \emph{같은 $s_{\phi,j}$,
같은 $n_j^{\eff}F$ scale} 을 비가역 \emph{heat-force} $\eta_j$ 에 적용한 형이다(rate-affinity 는 중심 $U_j$
기준, heat-force 는 국소 평형 $V_{\eq,j}$ 기준 — \S\ref{sec:ch4_eta_single} boundbox 의 구분 보존).

\emph{(5단계) entropy production 을 발열로(미시=거시).} 식~\eqref{eq:ch4_tr_entropy} 를 $\dot n_j$ 와
$\mathcal A_j^{\chem}$ 로 다시 쓰면 $\dot{\mathcal Q}_{\irr,j}^{\tr}=(Q_{j,\tot}/F)(\mathcal J_j^+-\mathcal J_j^-)
\cdot RT\ln(\mathcal J_j^+/\mathcal J_j^-)=\dot n_j\,\mathcal A_j^{\chem}$ 다($\dot n_j=(Q_{j,\tot}/F)(\dd\xi_j/\dd t)
=(Q_{j,\tot}/F)(\mathcal J_j^+-\mathcal J_j^-)$). 여기에 식~\eqref{eq:ch4_affinity_eta} 를 넣으면 $F$ 가 상쇄되어
\begin{equation}
\boxed{\;\dot{\mathcal Q}_{\irr,j}^{\tr}=\dot n_j\,\mathcal A_j^{\chem}
=\frac{Q_{j,\tot}}{F}\frac{\dd\xi_j}{\dd t}\cdot n_j^{\eff}F\,\eta_j
=n_j^{\eff}\,I_j\,\eta_j\;.}
\label{eq:ch4_micro_macro}
\end{equation}
즉 \textbf{미시 transition entropy production 이 거시 flux$\times$overpotential 형으로 정확히 환원}되며, ideal
width $w_j=RT/F$($n_j^{\eff}=1$)에서 $\dot{\mathcal Q}_{\irr,j}^{\tr}=I_j\eta_j$ 다. \textbf{부호(2법칙
재확인).} $\xi_j<\xi_{\eq,j}\Rightarrow$ (정의상) $V_{\eq,j}(\xi_j)<V_{n,\drive}\Rightarrow\eta_j>0$, 동시에
$\dd\xi_j/\dd t=k_j^{\mathrm{phys}}(\xi_{\eq,j}-\xi_j)>0\Rightarrow I_j>0$ — 둘이 같은 부호라 $n_j^{\eff}I_j\eta_j>0$
($n_j^{\eff}>0$). 반대로 $\xi_j>\xi_{\eq,j}$ 이면 둘 다 음이라 곱은 여전히 양; 평형서 $\eta_j\to0,\ I_j\to0$
동시 소멸. \textbf{이로써 Ch2(거시 flux$\times$overpotential)$\cdot$Ch3(미시 forward/backward rate, $n_j^{\eff}$)
$\cdot$Ch4(entropy production)가 한 표현 $n_j^{\eff}I_j\eta_j$ 로 닫힌다}(P3-6 미시=거시 정합 통과).
\textbf{일반/특수 위계.} \emph{Chapter 2 의 $\sum_j J_{n,j}F\eta_j=\sum_j I_j\eta_j$ 는 이 결과의
$n_j^{\eff}=1$(ideal-width) 특수해}이고, 일반 dissipation 은 $\sum_j n_j^{\eff}I_j\eta_j$ 다(\S\ref{sec:ch4_general})
— 두 장은 모순이 아니라 일반/특수 위계다(Chapter 3 의 $n_j^{\eff}=RT/(Fw_j)$ refinement 와 동일 맥락).
\end{verifybox}

\textbf{유도가 쓴 전제의 명시(★ 이중정의 차단 재확인).} 위 verifybox 는 (i) detailed balance(Ch3, \emph{내부}
$V_n$ 기준 ln 비), (ii) $n_j^{\eff}=RT/(Fw_j)$(Ch3), (iii) \textbf{기본형 $V_{n,\drive}=V_n$}
(3단계서 명시 전제)을 사용했다. (iii)이 없으면 detailed balance 가 쓴 $V_n$ 과 dissipation force 가 쓴
$V_{n,\drive}$ 가 서로 달라 $\mathcal A_j^{\chem}=n_j^{\eff}F\eta_j$ 가 닫히지 않는다 — 이 전제를 사용 시점에
박는 것이 Phase A CRITICAL($\eta$ 이중정의)의 해소다. 또 (iv) Chapter 3 의 keystone 한정(Level A $=$ Level B
의 detailed-balance 극한; $\xi_{\sstat,j}\to\xi_{\eq,j}$ 가 $C_j$ 의 $V_n$$\cdot$$\xi_j$-무의존 partition 에서만
성립)을 상속한다. 강구동 branch 비대칭(non-detailed-balance, Chapter 5)에서는 $\mathcal A_j^{\chem}$ 가 일반화된다.

\begin{boundbox}
\textbf{log singularity 는 floor 아닌 domain guard (\AL{49}).} $\mathcal J_j^\pm\to0$ 에서 $\ln(\mathcal J_j^+/\mathcal J_j^-)$
발산은 coarse-graining resolution 에 해당하는 작은 positive floor 로 처리한다 — 이는 heat fitting parameter 가
아니라 \emph{numerical domain guard} 다(GS-4 의 cap 금지와 구분: 물리 모델의 임계값이 아니라 수치 정의역 보호).
\textbf{게다가} 식~\eqref{eq:ch4_micro_macro} 형 $n_j^{\eff}I_j\eta_j$ 를 쓰면 $\mathcal J_j^\pm$ 가 평형 부근에서
작아도 $I_j=Q_{j,\tot}(\dd\xi_j/\dd t)\to0$ 으로 \emph{매끄럽게} 0 이 되어(그리고 $\eta_j\to0$) 발산 자체가
나타나지 않는다 — 즉 거시형은 미시형의 log 발산을 자동으로 정칙화한다. 따라서 실제 발열 계산은 거시형
$n_j^{\eff}I_j\eta_j$ 를 쓰고, 미시형 식~\eqref{eq:ch4_tr_entropy} 는 \emph{출처(entropy production)와 부호
($\ge0$) 의 근거}로 둔다(BOUNDED).
\end{boundbox}

\textbf{채택 조건(후보식의 검증).} 식~\eqref{eq:ch4_tr_entropy}--\eqref{eq:ch4_micro_macro} 를 실제 발열항으로
\emph{확정}하려면 (1) $r_j^\pm$ 가 실제 coarse-grained transition rate, (2) $\xi_j$ 가 progress coordinate,
(3) detailed balance $\xi_{\eq,j}$ 정합, (4) calorimetry/rest-relaxation 독립 검증이 필요하다. 따라서 본 식은
해석 가능한 물리식이되, 독립 검증 없이 \emph{모든} relaxation heat 를 이 항으로 단정하지 않는다(Chapter 1$\cdot$2
의 forward-only$\cdot$식별성 원칙 계승).

% ====================================================================
\section{비가역열 일반형 $\sum_j n_j^{\eff}I_j\eta_j\ge0$ 과 일반/특수 위계}\label{sec:ch4_general}
% ====================================================================
\textbf{물리.} 식~\eqref{eq:ch4_micro_macro} 의 per-transition 결과를 모든 동시 진행 transition 에 대해 합하면
음극의 \emph{총} 계면 반응 비가역 발열률이다(\AL{45}):
\begin{equation}
\boxed{\;\dot{\mathcal Q}_{\irr,n}=\sum_{j=1}^{N_p} n_j^{\eff}\,I_j\,\eta_j\ \ge\ 0\;,\qquad
n_j^{\eff}=\frac{RT}{F\,w_j(T)}\;.}
\label{eq:ch4_qirr_general}
\end{equation}
\textbf{$\ge0$ 의 무비약 확인(항별 + 합).} \S\ref{sec:ch4_tr} verifybox 에서 각 항 $n_j^{\eff}I_j\eta_j$ 는
$\xi_j\lessgtr\xi_{\eq,j}$ 양쪽에서 $I_j$ 와 $\eta_j$ 가 같은 부호라 항별로 $\ge0$ 이다(detailed-balance
기본형 전제). 따라서 비음항들의 합 $\sum_j n_j^{\eff}I_j\eta_j\ge0$ 이 2법칙으로 보장된다. \textbf{차원 검산.}
$n_j^{\eff}$(무차원)$\times I_j[\mathrm A]\times\eta_j[\mathrm V]=\mathrm{A\cdot V}=\mathrm W$ — 발열률 단위.
\textbf{부호 양정치 형.} 단일 우세 transition 극한에서 이를 전류$\cdot$과전압 크기 곱으로 축약하면
Chapter 2 의 $q_\irr=|I|\eta\ge0$($\eta=\sum_j(|J_{n,j}|F/|I|)|\eta_j|$) 형과 정합하며, ``$I_j\eta_j$ 단독''
(부호 미정렬)이 아니라 부호가 이미 정렬된 양정치 합임을 유지한다.

\begin{keybox}
\textbf{일반형$\leftrightarrow$Ch2 특수형 위계.} 식~\eqref{eq:ch4_qirr_general} 의 \emph{일반형}
$\sum_j n_j^{\eff}I_j\eta_j$ 와 Chapter 2 의 \emph{특수형} $\sum_j I_j\eta_j$ 의 관계는 \emph{모순이 아니라
일반/특수 위계}다:
\begin{equation}
\underbrace{\dot{\mathcal Q}_{\irr,n}=\sum_j n_j^{\eff}I_j\eta_j}_{\text{Ch4 일반형}}
\ \xrightarrow{\ n_j^{\eff}=RT/(Fw_j)\to1\ (w_j=RT/F,\ \text{ideal})\ }\
\underbrace{\sum_j I_j\eta_j}_{\text{Ch2 특수형}}.
\label{eq:ch4_hierarchy}
\end{equation}
Chapter 2 는 ideal-width($n_j^{\eff}=1$)를 기본으로 써서 $\sum_j I_j\eta_j$ 를 얻었고(그 boundbox 에서 일반형을
본 장으로 forward-ref), 본 장은 effective width($w_j\ne RT/F$, 따라서 $n_j^{\eff}\ne1$)를 허용하는 일반형을
닫는다. 그 일반형이 미시 transition entropy production(식~\eqref{eq:ch4_tr_entropy})에서 유도되므로, $n_j^{\eff}$
가중은 임의 피팅 인자가 아니라 \emph{detailed-balance 정합이 강제한} 물리 인자다(Chapter 3 \S$n_j^{\eff}$
정의장과 동일 출처).
\end{keybox}

\textbf{과전압의 물리 성분 분해(Ch2 계승).} 총 과전압 $\eta$(또는 transition별 $\eta_j$)는 charge-transfer
($\eta_\ct$, Chapter 3 의 effective barrier 활성화), ohmic($\eta_\ohm$, Chapter 1 의 $R_n$ 분극),
transport/concentration($\eta_\conc$) 기여의 합으로 분해되며 각각 $\ge0$ 이다. 단 이 성분들을 \emph{발열 저항}
으로 쓸 때는 Chapter 3 의 $R_n\ne R_\ct\ne R_{n,\eff}$ 구분을 유지한다(\S\ref{sec:ch4_transport}).

% ====================================================================
\section{가역 열 (1): 평형 전위 온도계수 (Bernardi)}\label{sec:ch4_revOCV}
% ====================================================================
\textbf{물리.} 가역 열은 평형 entropy 가 준다 — 비가역 dissipation 과 \emph{분리}되는 열역학 기원 항이다.
Chapter 2 가 전하 보존식에서 유도한 음극 평형 전위 온도계수를 그대로 받는다($Q_\cell,Q_{j,\tot}$ 기준 용량
고정, \AL{20})\textemdash 외부 lookup 이 아니라 \emph{전하 보존식의 평형 특수해 미분}이다(P3-2):
\begin{equation}
\left(\frac{\partial V_{n,\OCV}}{\partial T}\right)_q
=-\,\frac{\dfrac{\partial Q_\bg}{\partial T}+\displaystyle\sum_j Q_{j,\tot}\left(\dfrac{\partial\xi_{\eq,j}}{\partial T}\right)_{V_n}}
{\dfrac{\partial Q_\bg}{\partial V_n}+\displaystyle\sum_j Q_{j,\tot}\dfrac{\partial\xi_{\eq,j}}{\partial V_n}}.
\label{eq:ch4_ocv_dvdt}
\end{equation}
방전 기준 음극 reduced reversible heat 는(Chapter 2 식~의 Bernardi balance)
\begin{equation}
\boxed{\;\dot{\mathcal Q}_{\rev,n}^{\OCV}=s_{\rev,n}^{\conv}\,|I|\,T\left(\frac{\partial V_{n,\OCV}}{\partial T}\right)_q\;.}
\label{eq:ch4_rev_ocv}
\end{equation}
\textbf{차원 검산.} $|I|\,T\,(\partial V/\partial T)=\mathrm A\cdot\mathrm K\cdot(\mathrm{V/K})=\mathrm{A\,V}=\mathrm W$
— 가역 발열률 단위. $s_{\rev,n}^{\conv}\in\{+1,-1\}$ 는 음극 전위$\cdot$전류 방향$\cdot$heat-generation-positive
규약을 명시한 뒤 고정하는 bookkeeping factor 다(피팅 X).
\textbf{부호가변성.} $q_{\rev,j}=I_j\Pi_j$ ($\Pi_j=T\partial U_j/\partial T$, 식 (L6))는 $\partial U_j/\partial T$
와 전류 방향에 따라 흡열/발열 모두 가능하다 — 비가역 열($\ge0$)과 달리 가역 열은 부호가변이며, 그래서 둘을 한
경험항으로 섞지 않는다(Chapter 2 keystone 의 ``가역열$=$entropy(방향) / 비가역열$=$소산(속도)'' 분업).

\textbf{금지되는 연결(P3-1).}
\begin{equation}
\dot{\mathcal Q}_{\rev,n}\ \not\propto\ |I|\,T\,\frac{\dd V_{n,\app}}{\dd T}.
\label{eq:ch4_forbid}
\end{equation}
$V_{n,\app}=V_n+s_I|I|R_n$ 에는 분극$\cdot$kinetic lag$\cdot$transport loss$\cdot$hysteresis-like gap$\cdot R_n(T)$
가 섞이므로 그 경로 미분은 reversible entropy coefficient 가 아니다 — 이를 가역열로 쓰면 비가역 소산을 가역
entropy 로 오계상한다(Chapter 2 식~의 금지 연결 계승).

\begin{verifybox}
\textbf{Bernardi 가역열 $=$ 미시 reaction entropy 합 (Ch2 정합 재확인, 무비약).} Chapter 2 (식 (L6))의
$\partial U_j/\partial T=\Delta S_{\rxn,j}/(s_{\phi,j}F)$ 를 per-transition 가역열에 넣는다. transition $j$ 의
reaction rate 는 $\dot n_j=I_j/(s_{\phi,j}F)$ [mol/s](Faraday 법칙)이고, 그 entropic 열률은 $T\Delta S_{\rxn,j}\dot n_j$
이므로
\begin{equation}
q_{\rev,j}=T\,\Delta S_{\rxn,j}\,\frac{I_j}{s_{\phi,j}F}
=I_j\,T\,\frac{\partial U_j}{\partial T}=I_j\,\Pi_j,
\label{eq:ch4_qrev_micro}
\end{equation}
($s_{\phi,j}F$ 와 $1/(s_{\phi,j}F)$ 상쇄). 즉 \textbf{미시 reaction entropy $\Delta S_{\rxn,j}$ 의 합}
$\dot{\mathcal Q}_{\rev}=\sum_j I_j\Pi_j=\sum_j I_j T\partial U_j/\partial T$ 가 거시 Bernardi 형
$s_{\rev,n}^{\conv}|I|T(\partial V_{n,\OCV}/\partial T)_q$ 와 같은 entropy 정보로 정합한다(Chapter 2 의
double-counting 차단 정합식과 동일 제약). \textbf{비가역과의 대비.} 본 절 가역열은 \emph{reaction} entropy
$\Delta S_{\rxn,j}$(평형, $\partial U/\partial T$)에서, \S\ref{sec:ch4_tr} 비가역열은 transition entropy
production(non-equilibrium, $\eta_j$)에서 나온다 — 둘은 별개 entropy 이며(Chapter 2 \AL{24}: reaction entropy
$\ne$ activation entropy), 한 뿌리($U_j(T)$ 와 $V_{\eq,j}(\xi_j)$)에서 갈라진 가역/비가역 두 갈래다.
\end{verifybox}

% ====================================================================
\section{가역 열 (2): transition entropy basis 와 double counting 금지}\label{sec:ch4_revxi}
% ====================================================================
Chapter 1$\cdot$2 의 핵심 내부 상태변수는 $\xi_j$ 다. 따라서 상변이 transition 의 entropy 변화는 \emph{진행률
시간율} $\dd\xi_j/\dd t$ 를 통해서도 표현된다 — 식~\eqref{eq:ch4_rev_ocv} 의 OCV 계수 방식과 \emph{같은 entropy
정보}를 더 미시적으로 쓴 것이다(\AL{20}). $\Delta S_j(\equiv\Delta S_{\rxn,j})$ = $j$번째 effective transition
의 방전 방향 mol electron(또는 Li-equiv) 1 mol 당 entropy change 로 두면
\begin{equation}
\dot{\mathcal Q}_{\rev,\xi}=\sum_{j=1}^{N_p}s_{\rev,\xi,j}^{\conv}\,T\,\Delta S_j\,\frac{Q_{j,\tot}}{F}\,\frac{\dd\xi_j}{\dd t}.
\label{eq:ch4_rev_xi}
\end{equation}
\textbf{차원 검산.} $T[\mathrm K]\cdot\Delta S_j[\mathrm{J/(mol\,K)}]\cdot(Q_{j,\tot}/F)[\mathrm{mol}]\cdot
(\dd\xi_j/\dd t)[\mathrm{s^{-1}}]=\mathrm{J/s}=\mathrm W$ — 가역 발열률. 직접 입력은 $\dd\xi_j/\dd t$(휴지 구간서
정전류 좌표 변환 금지, \S\ref{sec:ch4_rest}).

\textbf{double counting 금지(택일 + 정합 제약; \AL{25}).} 식~\eqref{eq:ch4_rev_ocv} 와 식~\eqref{eq:ch4_rev_xi}
는 \emph{같은 평형 entropy 정보}를 서로 다른 수준에서 표현하므로, 둘을 독립 발열항으로 동시에 자유롭게 더하면
double counting 이다(평형 $V_{n,\OCV}$ 온도계수에 이미 transition/background entropy 가 반영). 따라서 택일한다:
\begin{enumerate}[label=\Alph*),leftmargin=2.0em]
\item \textbf{OCV 중심}: 전체 가역 열을 식~\eqref{eq:ch4_rev_ocv} 로 대표; $\Delta S_j$ 는 분해 해석/파라미터
  interpretation 용으로만 쓴다.
\item \textbf{Transition entropy 중심}: $\Delta S_j,h_\bg$ 로 분해하되, OCV $\partial V/\partial T$ 는 독립항이
  아니라 $\Delta S_j,h_\bg$ 를 묶는 \emph{equilibrium consistency condition}(Chapter 2 정합식)으로 둔다:
\end{enumerate}
\begin{equation}
\boxed{\;s_{\rev,n}^{\conv}|I|\,T\left(\frac{\partial V_{n,\OCV}}{\partial T}\right)_q
=s_\bg^{\conv}\,T\,h_\bg\,\dot Q_\bg^{\prog}
+\sum_j s_{\rev,\xi,j}^{\conv}\,T\,\Delta S_j\,\frac{Q_{j,\tot}}{F}\,\frac{\dd\xi_{\eq,j}}{\dd t}\;.}
\label{eq:ch4_consistency}
\end{equation}
\textbf{부호 규약 정합 주의.} 좌$\cdot$우변의 $s^{\conv}$($s_{\rev,n}^{\conv},s_\bg^{\conv},s_{\rev,\xi,j}^{\conv}$)
는 모두 \emph{동일한 heat-positive 규약}에서 유도돼야 본 제약이 부호 모순 없이 닫힌다(Chapter 2 의 정합 제약과
동일 — 서로 다른 bookkeeping 규약으로 독립 도입하면 같은 reaction entropy 가 좌우에서 반대 부호로 들어와 제약이
깨진다). \textbf{준평형 전제.} 좌변 OCV 계수는 \emph{평형}($\xi_j=\xi_{\eq,j}$) 미분이고 우변 가역열 항
(식~\eqref{eq:ch4_rev_xi})은 \emph{실제} 진행률 $\dd\xi_j/\dd t$ 를 쓰므로, 두 표현이 ``같은 entropy 정보''로
정합하는 것은 \emph{준평형 극한}($\xi_j\approx\xi_{\eq,j}$, 따라서 $\dd\xi_j/\dd t\approx\dd\xi_{\eq,j}/\dd t$)에
한정된다 — 가역열의 정의역 자체가 준평형이므로 이 한정은 모순이 아니라 적용범위 명시다.

% ====================================================================
\section{Background heat 와 coordinate shift}\label{sec:ch4_bg}
% ====================================================================
$Q_\bg$ 는 배경 누적 용량이지 entropy 자체가 아니다. background reversible heat 를 쓰려면 별도 background
entropy coefficient 가 필요하다(\AL{26}):
\begin{equation}
h_\bg(V_n,T)\equiv\left(\frac{\partial V_{\bg,\eq}}{\partial T}\right)_{\mathrm{bg}},\qquad
\dot{\mathcal Q}_{\rev,\bg}=s_\bg^{\conv}\,T\,h_\bg(V_n,T)\,\dot Q_\bg^{\prog}.
\label{eq:ch4_qrev_bg}
\end{equation}
\textbf{total derivative 와 progress 의 구분(무비약, 중요).} 전하 보존식 (L1) 을 시간 미분하면 $Q_\bg$ 가
$(V_n,T)$ 의 함수이므로 전미분 연쇄법칙으로 $\dot Q_\bg^{\mathrm{tot}}=(\partial Q_\bg/\partial V_n)\dd V_n/\dd t
+(\partial Q_\bg/\partial T)\dd T/\dd t$ 다. 이를 progress 와 coordinate shift 로 나눈다:
\begin{equation}
\dot Q_\bg^{\mathrm{tot}}=\dot Q_\bg^{\prog}+\dot Q_\bg^{\coord},\qquad
\dot Q_\bg^{\prog}=\frac{\partial Q_\bg}{\partial V_n}\frac{\dd V_n}{\dd t},\qquad
\dot Q_\bg^{\coord}=\frac{\partial Q_\bg}{\partial T}\frac{\dd T}{\dd t}.
\label{eq:ch4_qbg_split}
\end{equation}
여기서 \emph{$\dot Q_\bg^{\prog}$(전위 구동 progress)만} heat-generating 반응 진행으로 식~\eqref{eq:ch4_qrev_bg}
에 들어가고, $\dot Q_\bg^{\coord}$(온도 변화에 따른 capacity-coordinate shift)는 그대로 reaction rate 로 넣지
않는다 — 이 구분을 빠뜨리면 단순 온도 drift 가 가역 발열로 오계상된다(Chapter 2 의 bgsplit 계승). 등온 정전류
구간서는 실용적으로 $\dot Q_\bg^{\prog}\approx|I|-\sum_j Q_{j,\tot}(\dd\xi_j/\dd t)$ 로 쓸 수 있으나(전류 보존,
Chapter 3 식), \emph{비등온서는 $\dot Q_\bg^{\coord}$ 항 때문에 이 근사를 그대로 쓰지 않는다}.
\begin{equation}
\dot Q_\bg^{\prog}\ \approx\ |I|-\sum_j Q_{j,\tot}\frac{\dd\xi_j}{\dd t}\qquad(\text{등온 정전류 한정}).
\label{eq:ch4_qbg_iso}
\end{equation}

% ====================================================================
\section{Transport 및 물리적 병목 발열}\label{sec:ch4_transport}
% ====================================================================
\textbf{물리.} 강구동 제한을 arbitrary cap 이 아니라 \emph{물리적 병목}으로 두는 Chapter 1$\cdot$3 의 원칙을
발열에도 유지한다. 선형 near-equilibrium 에서 transport dissipation 은 $I^2R$ 형으로 근사되나, 일반적으로는
flux$\times$force 가 더 기본이다(\AL{47}). 선형 근사:
\begin{equation}
\dot{\mathcal Q}_{\transp,n}\approx I^2 R_{\transp,n}\qquad(\text{선형 near-equilibrium}),
\label{eq:ch4_transport_i2r}
\end{equation}
고율/nonlinear transport:
\begin{equation}
\boxed{\;\dot{\mathcal Q}_{\transp,n}=\sum_\ell J_\ell\,X_\ell\ \ge\ 0\;,}
\label{eq:ch4_transport_ff}
\end{equation}
$J_\ell$=물질 flux, $X_\ell$=그 공액 force(chemical-potential / electrolyte-potential / concentration
gradient). 이는 식~\eqref{eq:ch4_fluxforce} 의 transport 채널 specialization 이며 2법칙으로 $\ge0$ 이다.
\textbf{차원 검산.} $I^2R=\mathrm{A^2\cdot\Omega}=\mathrm{A^2\cdot V/A}=\mathrm{A\,V}=\mathrm W$; flux$\times$force
$J_\ell X_\ell$ 도 공액쌍이라 [W] — 정합.
\begin{boundbox}
\textbf{세 저항의 구분 유지(\AL{47}; Ch1$\cdot$3 계승).} Chapter 1 의 $R_n$ 은 apparent voltage-axis
coefficient($V_{n,\app}-V_n=s_I|I|R_n$)이므로 $I^2R$ heat resistance 로 \emph{직접} 쓰지 않는다:
\begin{equation}
R_n\ne R_{n,\eff},\qquad R_n\ne R_{\ct,n},\qquad R_{\ct,n}\ne R_{n,\eff}.
\label{eq:ch4_Rsep}
\end{equation}
(i) $R_n$ $=$ apparent voltage-axis shift, (ii) $R_{\ct,n}=RT/(FA_\eff j_{0,n})$ $=$ charge-transfer
small-signal(Chapter 3), (iii) $R_{n,\eff}$ $=$ heat-generation effective($\dot{\mathcal Q}_\irr=I^2R_{n,\eff}$
형, 선형 응답 한정). $R_{n,\eff}$ 사용은 independent physical/calorimetric constraint 가 있을 때만 허용한다
(BOUNDED; 동시 자유 피팅 금지 — Chapter 3 의 식별성 chain).
\end{boundbox}

% ====================================================================
\section{휴지 구간 relaxation heat (Ch2 후보의 확정 분해)}\label{sec:ch4_rest}
% ====================================================================
\textbf{물리.} 휴지 구간에서는 $I=0,\ \dd q/\dd t=0$ 이지만 내부 진행률은 (L2) 로 계속 완화된다:
$\dd\xi_j/\dd t=r_j^+(1-\xi_j)-r_j^-\xi_j\ne0$. 외부 전류에 의한 ohmic heat 는 사라지나(전류 $0$), 내부
자유에너지 감소와 entropy production 은 남는다. Chapter 2 가 \emph{후보}로만 남긴 휴지 relaxation heat 를
본 장에서 가역/비가역으로 \emph{분해}한다(\AL{48}):
\begin{equation}
\dot{\mathcal Q}_{\xi,\rlx}=\dot{\mathcal Q}_{\rev,\xi}^{\rlx}+\dot{\mathcal Q}_{\irr,\xi}^{\rlx}.
\label{eq:ch4_rest_split}
\end{equation}
\textbf{비가역 성분은 식~\eqref{eq:ch4_micro_macro} 가 그대로 적용된다.} 휴지 중에도 \emph{내부} transition
current $I_j=Q_{j,\tot}(\dd\xi_j/\dd t)\ne0$ 이므로(외부 $|I|=0$ 이라도 transition 들이 서로 다른 속도로
완화하며 내부 전류를 주고받음), 비가역 후보는
\begin{equation}
\dot{\mathcal Q}_{\irr,\xi}^{\rlx}=\sum_j n_j^{\eff}\,I_j\,\eta_j\ \ge\ 0\qquad(\text{외부 }|I|=0,\ I_j=Q_{j,\tot}\dot\xi_j\ne0),
\label{eq:ch4_rest_irr}
\end{equation}
이고(식~\eqref{eq:ch4_qirr_general} 와 동형, 단 $\sum_j I_j$ 가 외부 전류가 아니라 \emph{내부} 전류이며, 총 내부
전류 보존 $\sum_j I_j+I_\bg^{\prog}=0$ 은 \emph{별개} 진술이다). \textbf{이 $\ge0$ 의 근거는 전류 보존이 아니다}:
휴지서도 모든 transition 이 \emph{공통 무대} $V_n$ 을 공유하므로 각 항에서 $I_j\,(\propto\xi_{\eq,j}-\xi_j)$ 와
$\eta_j\,(=V_n-V_{\eq,j}(\xi_j))$ 가 같은 부호이고(위 미시=거시 verifybox 의 항별 부호 논증을 휴지구간에 그대로
재적용), 따라서 항별 $n_j^{\eff}I_j\eta_j\ge0$ 에서 합 $\ge0$ 이 따라온다. 가역 후보는 transition entropy
basis(식~\eqref{eq:ch4_rev_xi}, 단 휴지서 $\dd\xi_j/\dd t$ 직접 사용)다. \textbf{좌표 변환 금지.} 휴지/가변전류 구간에서는 $|I|/Q_\cell$ 이 시간 상수가
아니므로 정전류 좌표 변환 $\dd\xi_j/\dd t=(|I|/Q_\cell)\dd\xi_j/\dd q$ 를 쓰지 않고 $\dd\xi_j/\dd t$ 를 직접 쓴다.
\begin{flagbox}
\textbf{가역/비가역 분리 비율은 독립 검증 필요(\AL{48}).} 식~\eqref{eq:ch4_rest_split} 의 비가역 성분
(식~\eqref{eq:ch4_rest_irr})은 entropy production 으로 $\ge0$ 임이 보장되나, 가역 성분(transition entropy
exchange)과의 \emph{분리 비율}은 forward/backward rate 의 detailed-balance 정합 정도에 달려 있다. detailed
balance 가 정확히 성립하면 식~\eqref{eq:ch4_micro_macro} 가 비가역 성분을 닫지만, branch 비대칭(Chapter 5)에서는
$\mathcal A_j^{\chem}$ 가 일반화되어 분리 비율이 달라진다. 따라서 독립 calorimetry/rest-temperature response
없이 가역/비가역을 임의 분리하지 않는다(Chapter 2 가 후보로 남긴 이유의 확정 분해이되, \emph{정량} 비율은
검증 전제).
\end{flagbox}

% ====================================================================
\section{전기--열 coupling 계산 순서와 온도 되먹임}\label{sec:ch4_order}
% ====================================================================
온도는 Chapter 1--3 의 다음 항들에 되먹임된다: $Q_\bg(V_n,T)$, $\xi_{\eq,j}(V_n,T)$, $U_j(T)$, $w_j(T)$,
$r_j^\pm(\cdot,T)$, $k_j^{\mathrm{phys}}(\cdot,T)$, $R_n(q,T,|I|)$, $\rho_j(g;T)$, $V_{n,\OCV}(q,T)$. 따라서
비등온 발열 계산에서 $V_n,\xi_j,T$ 를 독립 후처리하지 않고 같은 시간 적분 루프에서 갱신한다(P3-3 순환 의존성).
계산 순서:
\begin{enumerate}[label=\arabic*),leftmargin=2.0em]
\item 외부 전류로 $q(t)$ 갱신: $\dd q/\dd t=|I|/Q_\cell$.
\item 현재 $q,T,\{\xi_j\}$ 에서 전하 보존식 (L1) 으로 $V_n$ root-find.
\item $V_n,T$ 에서 $\xi_{\eq,j}$, $\mathcal A_j$, $r_j^\pm$ 또는 $k_j^{\mathrm{phys}}$, 그리고
  $\eta_j=V_{n,\drive}-V_{\eq,j}(\xi_j)$(기본형 $V_{n,\drive}=V_n$) 계산.
\item $\dd\xi_j/\dd t$, $I_j$, $\dot n_j$ 계산.
\item 비가역 열: transition entropy production / 일반형 $\sum_j n_j^{\eff}I_j\eta_j$(식~\eqref{eq:ch4_qirr_general})
  / transport flux$\times$force(식~\eqref{eq:ch4_transport_ff}).
\item 가역 열: OCV coefficient(식~\eqref{eq:ch4_rev_ocv}) 또는 transition entropy basis(식~\eqref{eq:ch4_rev_xi})
  중 택일(정합 제약 식~\eqref{eq:ch4_consistency}).
\item background heat 사용 시 $\dot Q_\bg^{\prog}$ 와 $\dot Q_\bg^{\coord}$ 분리(식~\eqref{eq:ch4_qbg_split}).
\item 내부 열원 $\dot{\mathcal Q}_{\src,n}$(식~\eqref{eq:ch4_decomp}) 계산.
\item 열수지식(식~\eqref{eq:ch4_balance})으로 $T(t)$ 갱신.
\item 갱신 $T(t)$ 는 위 되먹임 항들에 반영(Chapter 2 순환).
\end{enumerate}
(수치 적분기$\cdot$root-find solver 구현과 수렴 판정은 본 이론 장 범위 밖이며 Chapter 6 로 분리한다; GS-4 의
이론식/피팅 실무 분리.)
\begin{boundbox}
\textbf{되먹임이 ICA tail 에 미치는 영향(Ch2 계승).} 비가역 발열 $\to T\uparrow\to k_j\uparrow$
($L=|I|/(Q_\cell k_j)\downarrow$, Chapter 1 spectrum shift) $\to$ 꼬리 단축; 동시에 ideal $w_j=RT/F\uparrow$ 로
평형 peak 폭 증가($n_j^{\eff}=RT/(Fw_j)$ 변화). 두 효과가 경쟁하므로 비등온 ICA tail 을 단일 부호로 단정하지
않는다(Chapter 1$\cdot$2 의 온도 tail 3계층 분해와 정합). 본 장 비가역 발열 $\sum_j n_j^{\eff}I_j\eta_j$ 가
이 되먹임의 열원 항이다.
\end{boundbox}

% ====================================================================
\section{식별성 및 필요한 실험 제약}\label{sec:ch4_ident}
% ====================================================================
물리 우선 모델에서도 식별성은 사라지지 않는다. 해결책은 임의 clipping 이 아니라 독립 실험과 물리적 prior 다.
다음 항들은 강하게 상관된다:
\begin{equation}
\eta_n\leftrightarrow\mathcal A_n\leftrightarrow r_j^\pm\leftrightarrow k_{j,\lim}\leftrightarrow R_{n,\eff}\leftrightarrow\Delta S_j\leftrightarrow h_\bg\leftrightarrow\rho_j(g;T).
\label{eq:ch4_ident}
\end{equation}
\begin{longtable}{p{0.36\textwidth} >{\raggedright\arraybackslash}p{0.52\textwidth}}
\toprule 실험/자료 & 제한 가능한 항 \\ \midrule \endhead
저율 OCV 온도 series & $(\partial V_{n,\OCV}/\partial T)_q$, $\Delta S_j$, $h_\bg$ \\
pulse/current interrupt & immediate polarization, 일부 $R_n$, 일부 $R_\ct$ \\
EIS/small-signal & $R_\ct$, transport time constant, film \\
calorimetry & $\dot{\mathcal Q}_{\irr}$, $\dot{\mathcal Q}_{\rev}$, heat scale \\
rest relaxation 온도 응답 & $\dot{\mathcal Q}_{\xi,\rlx}$ 후보(가역/비가역 분리) 검증 \\
rate series & rate limitation, $r_j^\pm$, $k_{j,\lim}$, $\rho_j(g;T)$ \\
\bottomrule
\end{longtable}
독립 자료 없이 동시 자유 결정 금지: (1) OCV coefficient vs transition entropy basis(double counting),
(2) $R_n,R_\ct,R_{n,\eff}$, (3) $r_j^\pm,k_{j,\lim},\rho_j(g;T)$, (4) $\Delta S_j,h_\bg$ 와 rest relaxation
heat, (5) transport heat vs hysteresis-like voltage loss.
\begin{practicebox}
초기 탐색서 heat residual / effective heat gain / capped heat correction / apparent $I^2R$ 가 피팅을
안정화할 수 있으나 \emph{물리 모델 기본식이 아니다}. 논문용 해석에서는 transport / charge-transfer / entropy
coefficient / relaxation entropy production / hysteresis loss 로 분해하거나 단순 empirical residual 로 명시하고
주 결론에서 제외한다(GS-4; solver/numerical 구현은 Chapter 6).
\end{practicebox}

% ====================================================================
\section{Chapter 4 의 가정 원장 (AL-40$\sim$49)}\label{sec:ch4_al}
% ====================================================================
본 장에서 사용한 가정을 통합 ledger(RB\_AL\_MASTER) Ch4 그룹(AL-40$\sim$49)으로 정리한다. AL-40,41 은 기등재분
(Foundation), AL-42$\sim$49 는 본 장 S2 에서 채운 신규 정의다(Phase A 적발 dangling AL-H 계열의 실제 정의 부여
포함: H1$\to$AL-40/42, H2$\to$AL-41/43). tier = GROUNDED/BOUNDED/FLAGGED.
{\footnotesize
\setlength{\tabcolsep}{3pt}
\sloppy
\begin{longtable}{>{\raggedright\arraybackslash}p{0.045\textwidth} >{\raggedright\arraybackslash}p{0.335\textwidth} >{\raggedright\arraybackslash}p{0.115\textwidth} >{\raggedright\arraybackslash}p{0.225\textwidth} >{\raggedright\arraybackslash}p{0.155\textwidth}}
\toprule AL\# & 가정 진술 & tier & 근거 cite (DOI/ISBN) & 사용 식 \\ \midrule \endhead
AL-40 & entropy production $T\dot S_\irr=\sum_\alpha J_\alpha X_\alpha\ge0$ (flux--force, 2법칙) & GROUNDED & degrootmazur1962(book), onsager1931, prigogine1967(book) (onsager DOI 10.1103/\brk PhysRev.37.405) & \eqref{eq:ch4_fluxforce} \\
AL-41 & transition-level entropy production $=$ network/master-equation 형 $(\mathcal J^+{-}\mathcal J^-)\ln(\mathcal J^+/\mathcal J^-)\ge0$ & GROUNDED & schnakenberg1976 (DOI 10.1103/\brk RevModPhys.48.571) & \eqref{eq:ch4_Jpm},\allowbreak\eqref{eq:ch4_tr_entropy} \\
AL-42 & 전기화학 entropy production $=J_n\mathcal A_n\ge0$; $I_n\eta_n$ 은 부호 convention 종속 magnitude($J_n\mathcal A_n$ 이 기본형); rate-affinity $\mathcal A_j$(중심 $U_j$) $\ne$ heat-force $\eta_j$(국소 평형) & GROUNDED & degrootmazur1962(book), prigogine1967(book), bazant2013 (DOI 10.1021/\brk ar300145c) & \eqref{eq:ch4_eta_def},\allowbreak\eqref{eq:ch4_JA} \\
AL-43 & transition chemical affinity $\mathcal A_j^{\chem}=RT\ln(\mathcal J_j^+/\mathcal J_j^-)$; 부호 보조정리 $(\mathcal J^+{-}\mathcal J^-)\ln(\mathcal J^+/\mathcal J^-)\ge0$ (두 양수 차×로그비) & GROUNDED & schnakenberg1976, prigogine1967(book) (DOI 10.1103/\brk RevModPhys.48.571) & \eqref{eq:ch4_tr_entropy},\allowbreak\eqref{eq:ch4_affinity_raw} \\
AL-44 & ★ 미시$=$거시: DB(AL-30)$+n_j^\eff{=}RT/(Fw_j)$(AL-34) $\Rightarrow\mathcal A_j^{\chem}{=}n_j^\eff F\eta_j$, $\dot{\mathcal Q}_{\irr,j}^{\tr}{=}n_j^\eff I_j\eta_j$ (★ $\eta_j{=}V_{n,\drive}{-}V_{\eq,j}$, 기본형 $V_{n,\drive}{=}V_n$ 전제 — CHARTER step2 이중정의 해소; $s_{\phi,j}$ 복원 step1) & GROUNDED & schnakenberg1976, bazant2013, degrootmazur1962(book) (DOI 10.1103/\brk RevModPhys.48.571; 10.1021/\brk ar300145c) & \eqref{eq:ch4_affinity_eta},\allowbreak\eqref{eq:ch4_micro_macro} \\
AL-45 & ★ 비가역열 일반형 $\dot{\mathcal Q}_\irr{=}\sum_j n_j^\eff I_j\eta_j\ge0$ (2법칙); Ch2 $\sum_j I_j\eta_j$ 는 $n^\eff{=}1$ ideal-width 특수해 (CHARTER step3 일반/특수 위계) & GROUNDED & degrootmazur1962(book), onsager1931, schnakenberg1976 (DOI 10.1103/\brk PhysRev.37.405; 10.1103/\brk RevModPhys.48.571) & \eqref{eq:ch4_qirr_general},\allowbreak\eqref{eq:ch4_hierarchy} \\
AL-46 & Ch2 3-분해 → Ch4 4-분해 정련: $\dot{\mathcal Q}_{\irr}^{\mathrm{Ch2}}\supset\dot{\mathcal Q}_{\irr}^{\mathrm{Ch4}}+\dot{\mathcal Q}_{\transp}$ (계면반응 EP 와 transport dissipation 분리; 이중계상 금지) & BOUNDED & bernardi1985, rao1997 (DOI 10.1149/\brk 1.2113792; 10.1149/\brk 1.1837884) & \eqref{eq:ch4_decomp},\allowbreak\eqref{eq:ch4_ch2refine} \\
AL-47 & transport/병목 발열 $=$ flux--force $\sum_\ell J_\ell X_\ell$ (선형 near-eq 서 $I^2R_\transp$); $R_n\ne R_{n,\eff}\ne R_\ct$ 유지($R_{n,\eff}$ 는 독립 calorimetric 제약 시만) & BOUNDED & degrootmazur1962(book), newman, bardfaulkner (ISBN 978-0-\brk 471-47756-3; 978-0-\brk 471-04372-0) & \eqref{eq:ch4_transport_ff},\allowbreak\eqref{eq:ch4_Rsep} \\
AL-48 & 휴지 relaxation heat $\dot{\mathcal Q}_{\xi,\rlx}{=}\dot{\mathcal Q}_{\rev,\xi}^\rlx{+}\dot{\mathcal Q}_{\irr,\xi}^\rlx$; 외부 $I{=}0$ 라도 내부 $I_j{\ne}0$ 의 dissipation $n_j^\eff I_j\eta_j$ 존재; 가역/비가역 정량분리는 독립 calorimetry 필요 (Ch2 AL-29 후보 확정) & BOUNDED & schnakenberg1976, degrootmazur1962(book) (DOI 10.1103/\brk RevModPhys.48.571) & \eqref{eq:ch4_rest_split},\allowbreak\eqref{eq:ch4_rest_irr} \\
AL-49 & log singularity($\mathcal J^\pm{\to}0$)는 coarse-graining resolution positive floor (numerical domain guard, heat fitting 아님); $n_j^\eff I_j\eta_j$ 형은 평형서 $I_j{\to}0$ 로 매끄럽게 0 (발산 없음) & BOUNDED & schnakenberg1976 (DOI 10.1103/\brk RevModPhys.48.571) & \eqref{eq:ch4_tr_entropy} boundbox \\
\bottomrule
\end{longtable}
}
\textbf{AGP(Assumption Grounding Policy).} 모든 본문 가정식은 AL-\# 를 인용한다. 본 장에 FLAGGED \emph{전용}
항목은 없으며(AL-48 의 가역/비가역 정량 분리 비율은 BOUNDED+검증 전제), entropy production$\cdot$flux--force$\cdot$
network thermo 는 GROUNDED, 정련/transport/식별성/domain-guard 는 BOUNDED(유효범위 병기)다. 임의
clip/cap/softplus/step 정의식 0(GS-4; log 발산은 식~\eqref{eq:ch4_tr_entropy} domain guard + 거시형 정칙화).
solver/numerical 구현 0(Ch.6 분리). \textbf{DOI 정확 귀속(★ Ch3 오귀속 교훈 계승):} degrootmazur1962$\cdot$
prigogine1967 $=$ 교과서(DOI 없음), onsager1931 $=$ 10.1103/PhysRev.37.405, schnakenberg1976 $=$
10.1103/RevModPhys.48.571 — Onsager DOI 를 degrootmazur 책에 붙이는 류의 오귀속 0건.

% ====================================================================
\section{Chapter 4 자체 검수표}\label{sec:ch4_selfcheck}
% ====================================================================
{\small
\begin{longtable}{>{\raggedright\arraybackslash}p{0.74\textwidth} >{\raggedright\arraybackslash}p{0.16\textwidth}}
\toprule 검수 항목 & 결과 \\ \midrule \endhead
Chapter 1--3(RB) 를 수정 없이 적층; 전하 보존식이 발열로 대체되지 않음(P3-6) & 통과(\S\ref{sec:ch4_inherit}) \\
발열 항을 fitting residual 아닌 물리적 heat source term 으로 정의 & 통과(\S\ref{sec:ch4_decomp}) \\
★ $\eta_j$ 기준 전위 단일화 $\eta_j{=}V_{n,\drive}{-}V_{\eq,j}$, 기본형 $V_{n,\drive}{=}V_n$ (CHARTER step2 이중정의 해소) & 통과(\S\ref{sec:ch4_eta_single}, keybox) \\
비가역 열을 flux--force / entropy production 에서 출발($J_n\mathcal A_n\ge0$ 기본형) & 통과(식~\eqref{eq:ch4_JA}) \\
$I\eta$ 를 부호 convention 종속 magnitude 와 구분(``$I\eta$ 단독'' 회피; $q_\irr{=}|I|\eta$ 양정치) & 통과(\S\ref{sec:ch4_irr},\ref{sec:ch4_general}) \\
transition EP 식에 $Q_{j,\tot}/F$ 포함, $(\mathcal J^+{-}\mathcal J^-)\ln(\cdot)\ge0$ 무비약 2법칙 증명 & 통과(식~\eqref{eq:ch4_tr_entropy}) \\
★ 미시 transition EP $=$ 거시 $n_j^{\eff}I_j\eta_j$ 정합(무비약 5단계; $s_{\phi,j}$ 복원 step1) & 통과(verifybox,식~\eqref{eq:ch4_micro_macro}) \\
★ 미시$=$거시 대입서 기본형 $V_{n,\drive}{=}V_n$ 전제 명시(이중정의 차단 재확인) & 통과(verifybox 3단계, 전제 명시 단락) \\
★ 비가역열 일반형 $\sum_j n_j^{\eff}I_j\eta_j\ge0$ + Ch2 특수형($n^\eff{=}1$) 위계(CHARTER step3) & 통과(식~\eqref{eq:ch4_qirr_general},\eqref{eq:ch4_hierarchy}) \\
ideal width 서 $\to I_j\eta_j$(Ch2), 평형서 $I_j{\to}0$ 로 log 발산 회피(domain guard) & 통과(boundbox,\AL{49}) \\
가역 열을 OCV coefficient 또는 transition entropy basis 에서 출발 & 통과(\S\ref{sec:ch4_revOCV},\ref{sec:ch4_revxi}) \\
★ Bernardi 가역열 $=$ 미시 reaction entropy 합 정합 재확인(reaction $\ne$ activation entropy) & 통과(verifybox \S\ref{sec:ch4_revOCV}) \\
$V_{n,\app}$ 온도 미분을 가역 열로 직접 쓰지 않음(P3-1) & 통과(식~\eqref{eq:ch4_forbid}) \\
OCV 방식과 entropy basis double counting 정합 제약(동일 heat-positive 규약) & 통과(식~\eqref{eq:ch4_consistency}) \\
$Q_\bg$ progress 와 coordinate shift 구분; 등온 근사 한정 & 통과(식~\eqref{eq:ch4_qbg_split},\eqref{eq:ch4_qbg_iso}) \\
$R_n,R_\ct,R_{n,\eff}$ 구분; transport heat 를 flux--force/제한된 $I^2R$ 로만 & 통과(식~\eqref{eq:ch4_Rsep}) \\
휴지 relaxation heat 가역/비가역 확정 분해; 외부 $I{=}0$ 라도 내부 $I_j{\ne}0$ (Ch2 후보 확정) & 통과(\S\ref{sec:ch4_rest}) \\
전기--열 coupling 계산 순서 + 온도 되먹임 명시(P3-3 순환) & 통과(\S\ref{sec:ch4_order}) \\
모든 식 차원$\cdot$부호$\cdot$2법칙$\cdot$극한 검산 명시 & 통과(\S\ref{sec:ch4_tr},\ref{sec:ch4_general} 등) \\
``자명/clearly/쉽게/obviously'' 본문 0건(G-noleap) & 통과 \\
임의 heat correction/capped heat 를 기본식에 넣지 않음(실무 팁 격하; GS-4) & 통과 \\
모든 가정 AL-40$\sim$49 인용; BOUNDED 유효범위 병기; DOI 정확 귀속(오귀속 0) & 통과(\S\ref{sec:ch4_al}) \\
Ch5/Ch6 이월 항목 분리(충전 branch 부호 → Ch5) & 통과(\S\ref{sec:ch4_transmit}) \\
\bottomrule
\end{longtable}
}

% ====================================================================
\section{Chapter 5 $\cdot$ Chapter 6 으로 전달되는 항목}\label{sec:ch4_transmit}
% ====================================================================
\begin{enumerate}[label=\arabic*),leftmargin=2.0em]
\item \textbf{Chapter 5}: 충전 branch $\dot{\mathcal Q}_{\rev}$ 부호 convention($s_{\phi,j}^b$ branch 부호
  \emph{유도}); branch-dependent entropy coefficient; hysteresis heat vs polarization heat 분리;
  charge/discharge asymmetric affinity/rate($C_j$ 의 detailed-balance 이탈 $\Rightarrow\xi_{\sstat,j}\ne\xi_{\eq,j}$);
  full-cell voltage gap heat 해석. 본 장 일반형 $\sum_j n_j^{\eff}I_j\eta_j$ 의 부호 양정치는 detailed-balance
  기본형 전제이며, branch 비대칭에서 $\mathcal A_j^{\chem}$ 의 일반화가 Chapter 5 과제다.
\item \textbf{Chapter 6}: $\dd\xi_j/\dd t$ 기반 heat-rate 계산; stiff kinetics$+$thermal ODE 결합 수치 해법;
  직접 $g$-분포 적분서 heat-rate 동시 계산; calorimetry 부재 시 발열 항을 피팅하지 않고 forward prediction
  으로만 두는 제한(Chapter 1$\cdot$2 의 forward-only 원칙).
\end{enumerate}

% ====================================================================
\section{Chapter 4 의 결론}\label{sec:ch4_concl}
% ====================================================================
첫째, 비가역 열은 임의 heat correction 이 아니라 flux$\times$force entropy production $T\dot S_\irr=\sum_\alpha
J_\alpha X_\alpha\ge0$(식~\eqref{eq:ch4_fluxforce}), 전기화학적으로 $T\dot S_{\irr,n}=J_n\mathcal A_n\ge0$
(식~\eqref{eq:ch4_JA})에서 출발한다. ``$I\eta$ 단독''(부호 미정렬)은 entropy production 의 일반 정의가 아니다.

둘째, \textbf{미시 transition entropy production $\dot{\mathcal Q}_{\irr,j}^{\tr}=(Q_{j,\tot}/F)RT(\mathcal J_j^+-\mathcal J_j^-)
\ln(\mathcal J_j^+/\mathcal J_j^-)\ge0$ 은 거시 $n_j^{\eff}I_j\eta_j$ 로 정확히 환원}된다(식~\eqref{eq:ch4_micro_macro}).
이 환원은 (i) Chapter 3 detailed balance, (ii) $n_j^{\eff}=RT/(Fw_j)$, (iii) \textbf{기본형 $V_{n,\drive}=V_n$}
(★ $\eta$ 기준 전위 단일화, CHARTER step2)을 전제로 닫히며, $s_{\phi,j}$ 부호 인자가 명시된다(★ CHARTER step1).
이로써 Ch2(거시 flux$\times$overpotential)$\cdot$Ch3(미시 forward/backward rate, $n_j^{\eff}$)$\cdot$Ch4
(entropy production)가 한 표현으로 닫히고, 평형서 $I_j\to0$ 로 발산 없이 0 이 된다.

셋째, \textbf{비가역 발열 일반형은 $\dot{\mathcal Q}_\irr=\sum_j n_j^{\eff}I_j\eta_j\ge0$}(식~\eqref{eq:ch4_qirr_general})
이며, \emph{Chapter 2 의 $\sum_j I_j\eta_j$ 는 그 ideal-width($n_j^{\eff}=1$) 특수해}다(★ CHARTER step3 일반/
특수 위계, 모순 아님). $n_j^{\eff}$ 가중은 임의 피팅 인자가 아니라 detailed-balance 정합이 강제한 물리 인자다.

넷째, 가역 열은 평형 전위 온도계수(Bernardi, 식~\eqref{eq:ch4_rev_ocv}) 또는 transition entropy basis
(식~\eqref{eq:ch4_rev_xi})에서 출발하며, 두 방식을 독립항으로 동시에 자유롭게 더하면 double counting 이다
(정합 제약 식~\eqref{eq:ch4_consistency}). Bernardi 가역열이 미시 reaction entropy $\Delta S_{\rxn,j}$ 의 합과
정합함을 재확인했고(verifybox \S\ref{sec:ch4_revOCV}), reaction entropy 는 activation entropy 와 별개다.

다섯째, Chapter 1 의 $R_n$ 은 apparent voltage-axis coefficient 이므로 heat-generation resistance 가 아니다
($R_n\ne R_\ct\ne R_{n,\eff}$, 식~\eqref{eq:ch4_Rsep}). transport/병목 발열은 flux$\times$force(식~\eqref{eq:ch4_transport_ff})
또는 독립 제약된 $I^2R_{n,\eff}$ 로만 쓰며, 강구동 rate limitation 을 soft cap 으로 절단하지 않는다(GS-4).

여섯째, 휴지 구간에서도 외부 $|I|=0$ 이지만 내부 transition flux $I_j\ne0$ 의 dissipation $\sum_j n_j^{\eff}I_j\eta_j\ge0$
이 남는다(식~\eqref{eq:ch4_rest_irr}; Chapter 2 가 후보로 남긴 relaxation heat 의 확정 분해). 본 장의
$\dot{\mathcal Q}_{\irr},\dot{\mathcal Q}_{\rev},\dot{\mathcal Q}_{\rlx},\dot{\mathcal Q}_{\transp}$ 는 Chapter 5
의 charge/discharge branch 및 hysteresis heat 해석으로 전달된다. 이로써 Chapter 1 의 ``열역학$=$방향(가역/
entropy), 동역학$=$속도(비가역/소산)'' 분업이 entropy-production 수준에서 \emph{정량}되되 앞 장과 충돌하지 않는다.

\section*{Chapter 5 로의 전달}
본 장이 확정한 발열 일반형($\dot{\mathcal Q}_\irr=\sum_j n_j^{\eff}I_j\eta_j\ge0$, transition entropy
production, Bernardi 가역열)과 상태변수 위에서: Chapter 5 는 $C_j$ 의 detailed-balance 이탈
($\xi_{\sstat,j}\ne\xi_{\eq,j}$)로 충방전 히스테리시스의 branch 의존 target 과 부호($s_{\phi,j}^b$)를 유도하고,
hysteresis heat 와 polarization heat 를 분리한다. Chapter 6 은 본 장 발열률의 수치 적분과 calorimetry 부재 시
forward-prediction 제한을 다룬다.

\begin{thebibliography}{99}
\bibitem{degrootmazur1962} S. R. de Groot, P. Mazur, \emph{Non-Equilibrium Thermodynamics},
  North-Holland (1962); Dover (1984). [flux--force entropy production; 교과서, DOI 없음]
\bibitem{onsager1931} L. Onsager, ``Reciprocal Relations in Irreversible Processes. I,'' \emph{Phys. Rev.}
  \textbf{37}, 405 (1931). DOI: 10.1103/\brk PhysRev.37.405.
\bibitem{prigogine1967} I. Prigogine, \emph{Introduction to Thermodynamics of Irreversible Processes},
  3rd ed., Interscience (1967). [entropy production; 교과서, DOI 없음]
\bibitem{schnakenberg1976} J. Schnakenberg, ``Network Theory of Microscopic and Macroscopic Behavior of
  Master Equation Systems,'' \emph{Rev. Mod. Phys.} \textbf{48}, 571 (1976). DOI: 10.1103/\brk RevModPhys.48.571.
\bibitem{bazant2013} M. Z. Bazant, ``Theory of Chemical Kinetics and Charge Transfer Based on
  Nonequilibrium Thermodynamics,'' \emph{Acc. Chem. Res.} \textbf{46}, 1144 (2013). DOI: 10.1021/\brk ar300145c.
\bibitem{bernardi1985} D. Bernardi, E. Pawlikowski, J. Newman, ``A General Energy Balance for Battery
  Systems,'' \emph{J. Electrochem. Soc.} \textbf{132}, 5 (1985). DOI: 10.1149/\brk 1.2113792.
\bibitem{rao1997} L. Rao, J. Newman, ``Heat-Generation Rate and General Energy Balance for Insertion
  Battery Systems,'' \emph{J. Electrochem. Soc.} \textbf{144}, 2697 (1997). DOI: 10.1149/\brk 1.1837884.
\bibitem{thomasnewman2003} K. E. Thomas, J. Newman, ``Thermal Modeling of Porous Insertion Electrodes,''
  \emph{J. Electrochem. Soc.} \textbf{150}, A176 (2003). DOI: 10.1149/\brk 1.1531194.
\bibitem{newman} J. Newman, K. E. Thomas-Alyea, \emph{Electrochemical Systems}, 3rd ed., Wiley (2004).
  ISBN 978-0-\brk 471-47756-3.
\bibitem{bardfaulkner} A. J. Bard, L. R. Faulkner, \emph{Electrochemical Methods}, 2nd ed., Wiley (2001).
  ISBN 978-0-\brk 471-04372-0.
\end{thebibliography}

\end{document}
